%!TeX root=main-DS1.tex
\chapter{Iskazni ra\v cun}

% \v Zelja da se formuli\v se pouzdana procedura zaklju\v civanja koja vodi ka objektivnim i pouzdanim stavovima
% vodila je jo\v s stare Grke (pre svega stoike i Aristotela) da po\v cnu da razmi\v sljaju o tome kako da formalizuju
% procese mi\v sljenja i zaklju\v civanja. Ista ideja je bila na umu i Lajbnicu kada je pri\v cao o
% \enquote{univerzalnom jeziku} i \enquote{univerzalnoj proceduri} odlu\v civanja istinitosti, kao i D\v zor\v zu Bulu
% kada je pisao svoju \emph{Matemati\v cku analizu logike} u kojoj je postavio osnove iskaznog ra\v cuna.

\emph{Iskaz} je re\v cenica \v ciju istinitost mo\v zemo \emph{nepobitno} i \emph{jednozna\v cno} da utvrdimo.
Tipi\v cni primeri iskaza su matemati\v cke tvrdnje kao \v sto su
\begin{center}
  \begin{tabular}{l}
    \textit{\enquote{Postoje prirodni brojevi $k$, $l$ i $m$ takvi da je $k^2 + l^2 = m^2$}} (ta\v can iskaz), i\\
    \textit{\enquote{Simetrale strana trougla se seku na opisanoj kru\v znici}} (neta\v can iskaz).
  \end{tabular}
\end{center}
Izjave iz svakodnevnog \v zivota \emph{naj\v ce\v s\'ce nisu iskazi} jer mogu biti podlo\v zni ukusu, shvatanjima i interpretacijama.
Za potrebe ovog kursa iskazi ne govore o svakodnevnim situacijama, ve\'c o apstraktnim fenomenima tipi\v cnim za matematiku i ra\v cunarstvo.

\smallskip

\noindent
\begin{minipage}{3.5in}
  \hspace*{1.5em}Nama \'ce posebno biti interesantni iskazi koje govore o stanju memorije tokom izvr\v savanja programa.
  Na primer, pretpostavimo da jedan mali program ima dve celobrojne promenljive, \verb`x` i \verb`y`, i jednu logi\v cku promenljivu \verb`p`,
  i pretpostavimo da je tokom izvr\v savanja programa u jednom trenutku stanje u memoriji ra\v cunara opisano u tabeli pored.
  Tada su slede\'ce tvrdnje iskazi:
  $$
    \mathtt{x} \le 6, \quad \mathtt{y} < 0, \quad \mathtt{p}
  $$
\end{minipage}\hfill
\begin{minipage}{1.5in}
  \begin{center}
  \begin{tabular}{|lr|}
    \multicolumn{2}{c}{\textsc{Memorija}}\\
    \hline
    \multicolumn{2}{|c|}{ }\\
    \cline{2-2}
    \verb`x`: & \multicolumn{1}{|r|}{7}\\
    \cline{2-2}
    \verb`y`: & \multicolumn{1}{|r|}{9}\\
    \cline{2-2}
    \verb`p`: & \multicolumn{1}{|r|}{\texttt{true}}\\
    \cline{2-2}
    \multicolumn{2}{|c|}{ }\\
    \hline
  \end{tabular}
  \end{center}
\end{minipage}\newline

\noindent
jer mo\v zemo nepobitno i jednozna\v cno da utvrdimo da su prva dva iskaza neta\v cni, dok je tre\'ci iskaz ta\v can.
U ovom kontekstu iskaz mo\v ze da bude neki logi\v cki izraz ili logi\v cka promenljiva.

\section{Iskazne formule}

U programiranju je \v cesto potrebno ispitati da li va\v ze dva uslova istovremeno pre nego \v sto se pre\dj e na
ra\v cunanje vrednosti nekog izraza. Naravno, postoji na\v cin da se od dva ili vi\v se
jednostavnih uslova napravi jedan slo\v zeni uslov: jednostavni uslovi
se pove\v zu logi\v ckim veznicima \emph{konjunkcije} $\land$,
\emph{disjunkcije} $\lor$, \emph{negacije} $\lnot$, \emph{implikacije} $\Rightarrow$ i \emph{ekvivalencije} $\Leftrightarrow$, recimo
ovako:
$$
  (x > 0 \;\;\land\;\; y < 0) \;\;\Rightarrow\;\; xy < 0.
$$
U savremenim programskim jezicima se uglavnom javljaju samo prva tri veznika, naj\v ce\v s\'ce u slede\'cem obliku:
\emph{konjunkcija} \verb"&&", \emph{disjunkcija} \verb"||" i \emph{negacija} \verb"!".
Evo nekoliko primera:

\begin{center}
    \begin{tabular}{@{}cc@{}}
    \hline
    Matematika & Programiranje \\
    \hline
    $z = 7$ & \verb"z == 7" \\
    $x < 5 \land y \ge 0$ & \verb"x < 5 && y >= 0" \\
    $\lnot(x = y \lor z > 0)$ & \verb"!(x == y || z > 0)" \\
    $\lnot\mathit{ok} \;\;\land\;\; \mathit{pos} \ge 0$ & \verb"!ok && pos >= 0" \\
    $(x < 2 \lor y \ne 4) \land (x > 3 \lor y > 3)$ &
    {\verb"(x < 2 || y != 4) && (x > 3 || y > 3)"}
    \end{tabular}
\end{center}

Kada analiziramo logi\v cka svojstva ovakvih slo\v zenih uslova obi\v cno se
prosti uslovi zamene slovima da bi se izrazi kojima manipuli\v semo pojednostavili.
Na primer, slo\v zeni uslov:
$$
  ( \, \underbrace{x < 2}_p \,\lor\, \underbrace{y \ne 4}_q \, ) \,\land\, (\, \underbrace{x > 3}_s \,\lor\, \underbrace{y > 3}_t \, )
$$
tako postaje \emph{iskazna formula}:
$$
  (p \lor q) \land (s \lor t).
$$
Slova $p$, $q$, $s$ i $t$ koja u\v cestvuju u gra\dj enju ove iskazne formule se zovu \emph{iskazna slova}.
Ona predstavljaju najjednostavnije iskaze \v cijim povezivanjem je nastao polazni slo\v zeni iskaz.

\begin{quote}
  \emph{Apstrahovanje} je proces u kome se neki aspekti fenomena zanemare kako bi se na\v sa pa\v znja fokusirala samo na one
  aspekte koji su klju\v cni za razmatranje. Za takva razmatranja onda ka\v zemo da su \emph{apstraktna}.
\end{quote}

U ovom delu kursa \'cemo razmatratrati uzajamne interakcije logi\v ckih pojmova bez ula\v zenja u
prirodu tih logi\v ckih pojmova. Tako dolazimo do slede\'ce \emph{apstraktne} definicije iskazne formule.


\ifBIOINF
      Neka je $p_1, \ldots, p_n$ kona\v can niz simbola koje zovemo \emph{iskazna slova}.
      \emph{Iskazna formula sa slovima $p_1, \ldots, p_n$}
      je kona\v can niz simbola koji se formira sli\v cno algebarskim izrazima:
	  jednostavnije izraze povezujemo u slo\v zenije \emph{binarnim veznicima}
	  $\land$, $\lor$, $\Rightarrow$, $\Leftrightarrow$, ili ih napadamo unarnim veznikom $\lnot$.
	  Pri tome koristimo zagrade da bismo naglasili redosled izvo\dj enja operacija.
\else
    \begin{DEF}
      Neka je $p_1, \ldots, p_n$ kona\v can niz simbola koje zovemo \emph{iskazna slova}.
      \emph{Iskazna formula sa slovima $p_1, \ldots, p_n$}
      je kona\v can niz simbola koji je formiran prema slede\'cim pravilima:
      \begin{enumerate}\itemsep -1mm
      \item
        simboli $\top$, $\bot$ i
        \emph{iskazna slova} $p_1$, $p_2$, \ldots, $p_n$ su iskazne formule;
      \item
        ako su $\phi$ i $\psi$ iskazne formule, onda su to i slede\'ci nizovi simbola:
        $$
          \lnot \phi, \quad (\phi \land \psi), \quad (\phi \lor \psi), \quad (\phi \Rightarrow \psi), \quad (\phi \Leftrightarrow \psi);
        $$
      \item
        svaka iskazna formula se dobija kona\v cnom primenom pravila (1) i (2).
      \end{enumerate}
      Zahtevom (3) se obezbe\dj uje da je iskazna formula kona\v can niz simbola.
    \end{DEF}
\fi

\begin{EX}
  Evo nekoliko iskaznih formula:
  \begin{itemize}
    \item $p_{19}$ je iskazna formula sa slovima $p_1$, \ldots, $p_{20}$ (primetimo da nismo u obavezi da upotrebimo sva iskazna slova u formuli);
    \item $(x \lor y)$ je iskazna formula sa slovima $x$, $y$;
    \item $((p \land \lnot p) \Rightarrow \bot)$ je iskazna formula sa slovom $p$;
    \item $(\lnot \lnot(u \lor v) \land (w \Rightarrow (v \Leftrightarrow u)))$ sa slovima $u$, $v$,$w$.
  \end{itemize}
\end{EX}

\begin{CONVENTION}
  Vidimo da iskazne formule mogu da budu prili\v cno komplikovani nizovi simbola.
  Zato se prilikom zapisivanja formula usvajaju slede\'ce konvencije:
  \begin{itemize}
  \item
    po\v sto se slova jasno vide u iskaznoj formuli ne\'cemo vi\v se govoriti \enquote{iskazna formula sa slovima \ldots},
    ve\'c samo \enquote{iskazna formula};
  \item
    nikada ne pi\v semo spolja\v snje zagrade kod formula; i
  \item
    veznici $\land$ i $\lor$ imaju vi\v si prioritet nego veznici~$\Rightarrow$ i~$\Leftrightarrow$.
  \end{itemize}
\end{CONVENTION}

\begin{EX}
  Uz ove konvencije formule iz prethodnog primera dobijaju slede\'ci (\v citkiji) oblik:
  $$
    \begin{array}{c@{\quad\text{ pi\v semo ovako: }\quad}c}
      (x \lor y) & x \lor y\\
      ((p \land \lnot p) \Rightarrow \bot) & p \land \lnot p \Rightarrow \bot \\
      (\lnot \lnot(u \lor v) \land (w \Rightarrow (v \Leftrightarrow u)))
      & \lnot \lnot(u \lor v) \land (w \Rightarrow (v \Leftrightarrow u))
    \end{array}
  $$
\end{EX}

Vrednosti logi\v ckih izraza u programiranju su obi\v cno ozna\v cene sa \verb`true` (ta\v cno) i \verb`false` (neta\v cno), dok ih
u matematici ozna\v cavamo sa $\top$ (ta\v cno) i $\bot$ (neta\v cno) i zovemo \emph{istinitosne vrednosti}.
Ako su poznate istinitosne vrednosti iskaznih slova koja u\v cestvuju u gra\dj enju iskazne formule
onda je lako odrediti istinitosnu vrednost cele formule
koriste\'ci \v cinjenicu da su logi\v cki veznici ujedno i logi\v cke operacije koje su definisane slede\'cim tablicama:
$$
  \begin{array}{ccccc}
    \begin{array}{|c|c|}
    \hline
        & \lnot \\
    \hline
    \bot & \top \\
    \top & \bot \\
    \hline
    \end{array}
  &
    \begin{array}{|c|cc|}
    \hline
    \land & \bot & \top \\
    \hline
    \bot & \bot & \bot \\
    \top & \bot & \top \\
    \hline
    \end{array}
  &
    \begin{array}{|c|cc|}
    \hline
    \lor  & \bot & \top \\
    \hline
    \bot & \bot & \top \\
    \top & \top & \top \\
    \hline
    \end{array}
  &
    \begin{array}{|c|cc|}
    \hline
    \Rightarrow  & \bot & \top \\
    \hline
    \bot & \top & \top \\
    \top & \bot & \top \\
    \hline
    \end{array}
  &
    \begin{array}{|c|cc|}
    \hline
    \Leftrightarrow  & \bot & \top \\
    \hline
    \bot & \top & \bot \\
    \top & \bot & \top \\
    \hline
    \end{array}
  \end{array}
$$

\medskip

\noindent
\begin{minipage}{3.5in}
  \begin{EX}
    Odrediti istinitosnu vrednost iskazne formule $((\lnot p \lor q) \land r) \Rightarrow (q \land r)$
    ako je istinitosna vrednost slova koja se javljaju u formuli data tabelom pored.
  \end{EX}
\end{minipage}\hfill
\begin{minipage}{1.5in}
  $$
    \begin{array}{ccc}
    p & q & r \\
    \hline
    \bot & \top & \bot
    \end{array}
  $$
\end{minipage}\par

\bigskip

\noindent
\begin{minipage}[t]{2in}
  \begin{proof}[Re\v senje]
    Ako u formuli slova zamenimo njihovim istinitosnim vrednostima i izvr\v simo nazna\v cene
    logi\v cke operacije dobijamo ra\v cun prikazan pored.
  \end{proof}
\end{minipage}\hfill
\begin{minipage}[t]{3.5in}
  \begin{center}
  \begin{tabular}[t]{l}
  $ ((\lnot \bot \lor \top) \land \bot) \Rightarrow (\top \land \bot)$\\
  $\;\;=\;\; ((\top \lor \top) \land \bot) \Rightarrow (\top \land \bot) $\\
  $\;\;=\;\; (\top  \land \bot) \Rightarrow (\top \land \bot) $\\
  $\;\;=\;\; \bot \Rightarrow (\top \land \bot) $\\
  $\;\;=\;\; \bot \Rightarrow \bot $\\
  $\;\;=\;\; \top$
  \end{tabular}
  \end{center}
\end{minipage}\par

\begin{DEF}
  Tabela koja svakom iskaznom slovu neke formule dodeljuje jednu od vrednosti
  $\top$ ili $\bot$ se zove \emph{valuacija}.
\end{DEF}

Istinitosna vrednost formule za datu valuaciju
se onda ra\v cuna pravolinijski, koriste\'ci tabelice koje defini\v su logi\v cke operacije.
U prethodnom primeru data nam je valuacija
  $$
    v : \begin{array}{ccc}
    p & q & r \\
    \hline
    \bot & \top & \bot
    \end{array}
  $$
\v Cinjenicu da u valuaciji $v$ slovo $p$ ima istinitosnu vrednost $\bot$ pi\v semo ovako:
$$
  v(p) = \bot,
$$
a \v cinjenicu da u valuaciji $v$ formula $\phi = ((\lnot p \lor q) \land r) \Rightarrow (q \land r)$
ima istinitosnu vrednost $\top$ pi\v semo ovako:
$$
  v(\phi) = \top.
$$

\section{Tautologije}

Iskazna formula mo\v ze u nekoj valuaciji da bude ta\v cna, a u nekoj neta\v cna. Me\dj utim,
postoje iskazne formule koje su uvek ta\v cne nezavisno od valuacije.
Na primer, takva je formula $\top$, ali i formule $p \lor \lnot p$ i $p \Rightarrow (q \Rightarrow p)$.
Takve formule se zovu tautologije.

\begin{DEF}
  Iskazna formula je \emph{tautologija} ako je ta\v cna za svaku valuaciju slova koja se u njoj javljaju.
  Izjavu \enquote{\textit{$\phi$ je tautologija}} skra\'ceno zapisujemo ovako:
  $$
    \vDash \phi
  $$
\end{DEF}

\begin{EX}
  Formula $\psi = ((\lnot p \lor q) \land r) \Rightarrow (q \land r)$ nije tautologija zato \v sto
  u slede\'coj valuaciji
  $
    v : \begin{array}{ccc}
    p & q & r \\
    \hline
    \bot & \bot & \top
    \end{array}
  $
  za istinitosnu vrednost formule dobijamo:
  \begin{align*}
    v(\psi) &= ((\lnot \bot \lor \bot) \land \top) \Rightarrow (\bot \land \top) = ((\top \lor \bot) \land \top) \Rightarrow (\bot \land \top) \\
    &= (\top \land \top) \Rightarrow (\bot \land \top) = \top \Rightarrow (\bot \land \top) = \top \Rightarrow \bot = \bot
  \end{align*}
  Dakle, ako postoji \emph{bar jedna} valuacija za koju formula nije ta\v cna, ta formula \emph{nije} tautologija.
\end{EX}

\begin{EX}
  Dokazati da je formula $\phi = (p \Rightarrow q) \Leftrightarrow (\lnot q \Rightarrow \lnot p)$ tautologija.
\end{EX}
\begin{proof}[Re\v senje.]
  Ako formula ima svega nekoliko slova, najbr\v zi na\v cin da se proveri da li je ona tautologija je
  da izra\v cunamo njenu vrednost za sve mogu\'ce valuacije, \v sto je ura\dj eno u Tabeli~\ref{ft03.tab.primer-taut-0}.
  Ovakav na\v cin provere da li je formula tautologija ima smisla za formule u kojima se javlja relativno malo slova (do \v cetiri)
  zato \v sto za formulu sa $n$ slova ovakva tabela, pored zaglavlja, ima jo\v s $2^n$ redova.
  Na primer, za formulu koja ima 10 iskaznih slova ovakva tabela ima zaglavlje i jo\v s 1024 reda!
\end{proof}

\begin{table}[!h]
    $$
      \begin{array}{|cccc|c|c||c|}
      \hline
      p & q & p \Rightarrow q & \lnot q \Rightarrow \lnot p & \phi \\
      \hline
      \bot & \bot & \top & \top & \top \\
      \bot & \top & \top & \top & \top \\
      \top & \bot & \bot & \bot & \top \\
      \top & \top & \top & \top & \top \\
      \hline
      \end{array}
    $$
    \caption{Dokaz da je formula $\phi = (p \Rightarrow q) \Leftrightarrow (\lnot q \Rightarrow \lnot p)$ tautologija}
    \label{ft03.tab.primer-taut-0}
\end{table}

Nekada se de\v sava da nije lako na\'ci direktan dokaz nekog stava. Tada se \v cesto pribegava lukavstvu
koje se zove \emph{svo\dj enje na kontradikciju} (lat.~\emph{reductio ad absurdum}) i zasniva se na tautologiji
$$
  \vDash p \Leftrightarrow (\lnot p \Rightarrow \bot)
$$
koja ka\v ze da je logi\v cki svejedno da li \'cemo dokazivati iskaz $p$ ili \'cemo dokazati da
pretpostavka da ne va\v zi $p$ nu\v zno vodi do neta\v cnog iskaza.


\begin{EX}
  Dokazati da je formula $\phi = (x \Rightarrow y) \land (z \Rightarrow w) \Rightarrow (x \lor z \Rightarrow y \lor w)$
  tautologija.
\end{EX}
\begin{proof}[Re\v senje]
%  Jedan na\v cin da to poka\v zemo je da proverimo sve mogu\'ce valuacije sa \v cetiri slova, \v sto je ura\dj eno
%  u Tabeli~\ref{ft03.tab.primer-taut}.
%  Ovakav na\v cin provere da li je formula tautologija ima smisla za formule u kojima se javlja relativno malo slova (tri-\v cetiri)
%  zato \v sto za formulu sa $n$ slova ovakva tabela ima $2^n$ redova. (Na primer, za formulu koja ima 10 slova ovakva tabela ima
%  1024 reda!)
%  \begin{table}[!h]
%    $$
%      \begin{array}{|cccc|c|c||c|}
%      \hline
%      p & q & r & s & (p \Rightarrow q) \land (r \Rightarrow s) & p \lor r \Rightarrow q \lor s & \phi \\
%      \hline
%      \bot & \bot & \bot & \bot & \top & \top & \top \\
%      \bot & \bot & \bot & \top & \top & \top & \top \\
%      \bot & \bot & \top & \bot & \bot & \bot & \top \\
%      \bot & \bot & \top & \top & \top & \top & \top \\
%      \bot & \top & \bot & \bot & \top & \top & \top \\
%      \bot & \top & \bot & \top & \top & \top & \top \\
%      \bot & \top & \top & \bot & \bot & \top & \top \\
%      \bot & \top & \top & \top & \top & \top & \top \\
%      \top & \bot & \bot & \bot & \bot & \bot & \top \\
%      \top & \bot & \bot & \top & \bot & \top & \top \\
%      \top & \bot & \top & \bot & \bot & \bot & \top \\
%      \top & \bot & \top & \top & \bot & \top & \top \\
%      \top & \top & \bot & \bot & \top & \top & \top \\
%      \top & \top & \bot & \top & \top & \top & \top \\
%      \top & \top & \top & \bot & \bot & \top & \top \\
%      \top & \top & \top & \top & \top & \top & \top \\
%      \hline
%      \end{array}
%    $$
%    \caption{Dokaz da je formula $\phi = (p \Rightarrow q) \land (r \Rightarrow s) \Rightarrow (p \lor r \Rightarrow q \lor s)$
%    tautologija}
%    \label{ft03.tab.primer-taut}
%  \end{table}
%  
  Jedan na\v cin da to poka\v zemo je da proverimo sve mogu\'ce valuacije sa \v cetiri slova, kojih ima~16.
  Umesto toga, pokaza\'cemo drugi na\v cin da se proveri da li je formula tautologija.

  Da bismo dokazali da je formula $\phi$ tautologija treba da doka\v zemo da je ona ta\v cna za svaku mogu\'cu valuaciju.
  Dakle, treba da doka\v zemo iskaz:
  $$
    p = \text{\enquote{Za svaku valuaciju formula $\phi$ je ta\v cna.}}
  $$
  Primeni\'cemo strategiju dokazivanja svo\dj enjem na kontradikciju, prema kojoj je iskaz $p$ logi\v cki
  ekvivalentan iskazu $\lnot p \Rightarrow \bot$. Drugim re\v cima, pretpostavi\'cemo:
  $$
    \lnot p = \text{\enquote{Nije za svaku valuaciju formula $\phi$ ta\v cna}}
  $$
  pa \'cemo polaze\'ci od toga poku\v sati da dedukujemo neta\v can iskaz.
  
  Pretpostavimo da je $\lnot p$ ta\v cno, odnosno da nije za svaku valuaciju formula $\phi$ ta\v cna.
  To zna\v ci da postoji neka valuacija u kojoj je $\phi$ neta\v cna. Po\v sto je formula $\phi$
  oblika $\ldots \Rightarrow \ldots$, valuacija mora biti takva da je vrednost formule levo od znaka implikacije~$\top$,
  dok je vrednost formule desno od znaka implikacije~$\bot$. Dakle,
  $$
    v : \begin{array}{c@{\qquad}c}
      (x \Rightarrow y) \land (z \Rightarrow w) & x \lor z \Rightarrow y \lor w \\
      \hline
      \top & \bot
    \end{array}
  $$
  Konjunkcija dve formule je ta\v cna samo kada su obe formule ta\v cne, pa zaklju\v cujemo da je
  $$
    v : \begin{array}{c@{\qquad}c@{\qquad}c}
      x \Rightarrow y & z \Rightarrow w & x \lor z \Rightarrow y \lor w \\
      \hline
      \top & \top & \bot
    \end{array}
  $$
  S druge strane, implikacija je neta\v cna samo kada je formula levo od znaka implikacije ta\v cna, a formula desno od
  znaka implikacije neta\v cna, pa dobijamo:
  $$
    v : \begin{array}{c@{\qquad}c@{\qquad}c@{\qquad}c}
      x \Rightarrow y & z \Rightarrow w & x \lor z & y \lor w \\
      \hline
      \top & \top & \top & \bot
    \end{array}
  $$
  Disjunkcija je neta\v cna samo kada su oba disjunkta neta\v cni:
  $$
    v : \begin{array}{c@{\qquad}c@{\qquad}c@{\qquad}c@{\qquad}c}
      x \Rightarrow y & z \Rightarrow w & x \lor z & y & w \\
      \hline
      \top & \top & \top & \bot & \bot
    \end{array}
  $$
  Kako $v(x \Rightarrow y) = \top$, a $v(y) = \bot$, zaklju\v cujemo da mora biti $v(x) = \bot$:
  $$
    v : \begin{array}{c@{\qquad}c@{\qquad}c@{\qquad}c@{\qquad}c}
      x & z \Rightarrow w & x \lor z & y & w\\
      \hline
      \bot & \top & \top & \bot & \bot
    \end{array}
  $$
  Analogno, zbog $v(z \Rightarrow w) = \top$ i $v(w) = \bot$ mora biti $v(z) = \bot$:
  $$
    v : \begin{array}{c@{\qquad}c@{\qquad}c@{\qquad}c@{\qquad}c}
      x & z & x \lor z & y & w\\
      \hline
      \bot & \bot & \top & \bot & \bot
    \end{array}
  $$
  Tako smo dobili da je $\top = v(x \lor z) = v(x) \lor v(z) = \bot \lor \bot = \bot$, \v sto je neta\v can iskaz.
  Dakle, pretpostavka da postoji valuacija u kojoj polazna formula nije tautologija vodi u kontradikciju,
  \v sto zna\v ci da je ta pretpostavka neodr\v ziva. Time je dokazano da je $\phi$ tautologija.
\end{proof}

\begin{EX}
  Dokazati da je slede\'ca formula tautologija za svako $n \ge 1$:
  $$
    p_1 \Rightarrow (p_2 \Rightarrow (p_3 \Rightarrow ( \ldots \Rightarrow (p_n \Rightarrow p_1) \ldots )))
  $$
\end{EX}
\begin{proof}[Re\v senje]
  Pretpostavimo i ovaj put da formula nije tautologija. Onda postoji valuacija u kojoj je formula neta\v cna.
  Kako formula ima oblik $p_1 \Rightarrow ( \ldots )$, svaka valuacija u kojoj je ona bude neta\v cna izgleda ovako:
  $$
    v: \begin{array}{c@{\qquad}c@{\qquad}c@{\qquad}c@{\qquad}c}
      p_1 & p_2 \Rightarrow (p_3 \Rightarrow ( \ldots \Rightarrow (p_n \Rightarrow p_1) \ldots ))\\
      \hline
      \top & \bot 
    \end{array}
  $$
  Na isti na\v cin dalje zaklju\v cujemo da mora biti:
  $$
    v : \begin{array}{c@{\qquad}c@{\qquad}c@{\qquad}c@{\qquad}c}
      p_1 & p_2 & p_3 \Rightarrow ( \ldots \Rightarrow (p_n \Rightarrow p_1) \ldots )\\
      \hline
      \top & \top & \bot 
    \end{array}
  $$
  i tako dalje. Na kraju dobijamo da mora biti:
  $$
    v : \begin{array}{c@{\qquad}c@{\qquad}c@{\qquad}c@{\qquad}c}
      p_1 & p_2 & p_3 & \ldots & p_n \Rightarrow p_1\\
      \hline
      \top & \top & \top & \ldots & \bot 
    \end{array}
  $$
  i na kraju:
  $$
    v : \begin{array}{c@{\qquad}c@{\qquad}c@{\qquad}c@{\qquad}c@{\qquad}c}
      p_1 & p_2 & p_3 & \ldots & p_n & p_1\\
      \hline
      \top & \top & \top & \ldots & \top & \bot 
    \end{array}
  $$
  Tako smo dobili da je $\top = v(p_1) = \bot$, \v sto je neta\v can iskaz.
  Pretpostavka da postoji valuacija u kojoj polazna formula nije tautologija vodi u kontradikciju,
  \v sto zna\v ci da je ta pretpostavka neodr\v ziva. Time je dokazano da je $\phi$ tautologija.
\end{proof}

\begin{EX}
  Dokazati da je slede\'ca formula tautologija za sve $n \ge 2$:
  $$
    p \land (q_1 \lor \ldots \lor q_n) \Leftrightarrow (p \land q_1) \lor \ldots \lor (p \land q_n).
  $$
\end{EX}
\begin{proof}[Re\v senje]
  U ovom primeru \'cemo pokazati jo\v s jednu tehniku dokazivanja da je formula tautologija koja se zasniva na diskusiji
  po vrednosti iskaznog slova. Da bismo dokazali da je formula tautologija treba da doka\v zemo da je ta\v cna za sve valuacije.
  U proizvoljnoj valuaciji $v$ promenljiva $p$ mo\v ze imati vrednost ili $\top$ ili $\bot$.
  
  Slu\v caj 1$^\circ$: $v(p) = \top$. Formula tada postaje:
  $$
    \top \land (q_1 \lor \ldots \lor q_n) \Leftrightarrow (\top \land q_1) \lor \ldots \lor (\top \land q_n),
  $$
  odnosno,
  $$
    q_1 \lor \ldots \lor q_n \Leftrightarrow q_1 \lor \ldots \lor q_n,
  $$
  (zato \v sto je $v(\top \land u) = v(u)$), \v sto je o\v cito ta\v cno.

  Slu\v caj 2$^\circ$: $v(p) = \bot$. Formula tada postaje:
  $$
    \bot \land (q_1 \lor \ldots \lor q_n) \Leftrightarrow (\bot \land q_1) \lor \ldots \lor (\bot \land q_n),
  $$
  odnosno,
  $$
    \bot \Leftrightarrow \bot \lor \ldots \lor \bot,
  $$
  (zato \v sto je $v(\bot \land u) = \bot$), \v sto je o\v cito ta\v cno.
\end{proof}

Anti\v cki, a kasnije i srednjevekovni mislioci, imali su sna\v znu potrebu da formalizuju
\enquote{zakone pravilnog mi\v sljenja}. Tako su, kao va\v zni zakoni mi\v sljenja,
uo\v ceni slede\'ci \enquote{principi} koje sada navodimo u modernoj interpretaciji sa odgovaraju\'cim imenima
koja su skovana u srednjem veku. Prve dve formule koje navodimo predstavljaju dva od tri \enquote{klasi\v cna zakona mi\v sljenja}:

\medskip

\begin{tabular}{ll}
  $\vDash p \lor \lnot p$         & Zakon isklju\v cenja tre\'ceg (\textit{Tertium non datur})\\
  $\vDash \lnot(p \land \lnot p)$ & Princip kontradikcije (\textit{Principium contradictionis})\\
\end{tabular}

\medskip

\noindent
(tre\'ci \enquote{klasi\v cni zakon mi\v sljenja} glasi \enquote{sve je jednako sebi} i danas ga zapisujemo ovako: $x = x$.)
Evo i nekoliko \enquote{klasi\v cnih zakona zaklju\v civanja}:

\medskip

\begin{tabular}{ll}
  $\vDash p \land (p \Rightarrow q) \Rightarrow q$  & (\textit{Modus ponens})\\
  $\vDash (p \Rightarrow q) \land \lnot q \Rightarrow \lnot p$  & (\textit{Modus tollens})\\
  $\vDash (p \Rightarrow q) \Leftrightarrow (\lnot q \Rightarrow \lnot p)$  & Kontrapozicija\\
  $\vDash (p \Rightarrow q) \Leftrightarrow (p \land \lnot q \Rightarrow \bot)$  & Svo\dj enje na kontradikciju\\
                                                                                 & (\textit{Reductio ad absurdum})\\
  $\vDash (p \lor q) \land \lnot q \Rightarrow p$  & Disjunktivni silogizam\\
  $\vDash (p \Rightarrow q) \land (q \Rightarrow r) \Rightarrow (p \Rightarrow r)$  & Hipoteti\v cki silogizam\\
  $\vDash (p \Rightarrow q) \land (r \Rightarrow s) \Rightarrow (p \lor r) \Rightarrow (q \lor s)$  & Konstruktivna dilema\\
  $\vDash \bot \Rightarrow q$ & Iz la\v zi --- bilo \v sta\\
                                         & (\textit{Ex falso qoudlibet})
\end{tabular}

\medskip

Dugo vremena nije bilo jasno kako razdvojiti formule koje predstavljaju \enquote{zakone mi\v sljenja} od formula koje
to nisu. Predlog re\v senja ove dileme predstavlja okosnicu knjige \textit{Matemati\v cka analiza logike}
britanskog matemati\v cara D\v zord\v za Bula iz~1847.\ godine gde se predla\v ze da
\enquote{zakone mi\v sljenja} opisuju one i samo one formule koje su ta\v cne u svakoj valuaciji
(dakle, tautologije).

U svojoj detaljnoj analizi algebarskih svojstava iskaznog ra\v cuna D\v zord\v z Bul koristi i neke druge tautologije
koje ne li\v ce toliko na \enquote{zakone mi\v sljenja} koliko na \enquote{algebarska pravila} (odatle i poti\v ce ime
\enquote{iskazni \emph{ra\v cun}}), Tabela~\ref{ft03.tab.tautologije}.

\begin{table}[!h]
\centering
\begin{tabular}{ll}
    $\vDash \lnot \lnot p \Leftrightarrow p$ & Involutivnost negacije\\
    $\vDash p \lor \lnot p \Leftrightarrow \top$ & Zakon isklju\v cenja tre\'ceg\\
    $\vDash p \land \lnot p \Leftrightarrow \bot$ & Princip kontradikcije\\
    \hline
    $\vDash p \land \top \Leftrightarrow p$ & Neutralni element za $\land$\\
    $\vDash p \lor \bot \Leftrightarrow p$ & \hphantom{Neutralni element}\llap{\dots} za $\lor$\\
    \hline
    $\vDash p \land \bot \Leftrightarrow \bot$ & Apsorptivni element za $\land$\\
    $\vDash p \lor \top \Leftrightarrow \top$ & \hphantom{Apsorptivni element}\llap{\dots} za $\lor$\\
    \hline
    $\vDash p \land p \Leftrightarrow p$ & Idempotentnost konjunkcije\\
    $\vDash p \lor p \Leftrightarrow p$ & \hphantom{Idempotentnost}\llap{\dots} disjunkcije\\
    \hline
    $\vDash p \land q \Leftrightarrow q \land p$ & Komutativnost konjunkcije\\
    $\vDash p \lor q \Leftrightarrow q \lor p$ & \hphantom{Komutativnost}\llap{\dots} disjunkcije\\
    $\vDash (p \Leftrightarrow q) \Leftrightarrow (q \Leftrightarrow p)$ & \hphantom{Komutativnost}\llap{\dots} ekvivalencije\\
    \hline
    $\vDash (p \land q) \land r \Leftrightarrow p \land (q \land r)$ & Asocijativnost konjunkcije\\
    $\vDash (p \lor q) \lor r \Leftrightarrow p \lor (q \lor r)$ & \hphantom{Asocijativnost}\llap{\dots} disjunkcije\\
    \hline
    $\vDash p \land (q \lor r) \Leftrightarrow (p \land q) \lor (p \land r)$ & Distributivnost $\land$ prema $\lor$\\
    $\vDash p \lor (q \land r) \Leftrightarrow (p \lor q) \land (p \lor r)$ & \hphantom{Distributivnost}\llap{\dots} $\lor$ prema $\land$\\
    $\vDash (p \Rightarrow (q \Rightarrow r)) \Leftrightarrow ((p \Rightarrow q) \Rightarrow (p \Rightarrow r))$ & \hphantom{Distributivnost}\llap{\dots} $\Rightarrow$ prema sebi\\
    \hline
    $\vDash \lnot(p \lor q) \Leftrightarrow \lnot p \land \lnot q$ & De Morganovi zakoni \\
    $\vDash \lnot(p \land q) \Leftrightarrow \lnot p \lor \lnot q$ & \hphantom{De Morganovi zakoni}\llap{\dots}\\
    \hline
    $\vDash (p \Rightarrow q) \Leftrightarrow \lnot p \lor q$ & Zakon materijalne implikacije\\
    $\vDash (p \Leftrightarrow q) \Leftrightarrow (p \Rightarrow q) \land (q \Rightarrow p)$ & \hphantom{Zakon materijalne}\llap{\dots} ekvivalencije\\
\end{tabular}
\caption{\enquote{Algebarska pravila} iskaznog ra\v cuna}
\label{ft03.tab.tautologije}
\end{table}


\begin{CONVENTION}\label{ft03.convention.assoc}
  Asocijativnost konjunkcije i disjunkcije nam omogu\'cuju da jo\v s malo pojednostavimo pisanje formula. Tako \'cemo:
  umesto $p \land (q \land r)$ pisati kra\'ce $p \land q \land r$, i
  umesto $p \lor (q \lor r)$ pisati kra\'ce $p \lor q \lor r$.
\end{CONVENTION}



\section{Tautologije kao transformaciona pravila}

U ovom odeljku \v zelimo da poka\v zemo kako se navedene tautologije koriste kao \enquote{transformaciona pravila}.
Ideja je da napravimo alat koji \'ce nam omogu\'citi da transformi\v semo formule u nove formule koje mo\v zda
imaju jednostavniji ili \enquote{lep\v si} oblik, ali tako da se pri transformacijama logi\v cka vredost formule ne menja.

Da se podsetimo: iskazna formula je \emph{niz simbola (string)} formiran prema odgovaraju\'cim pravilima. Zbog toga
formule $\lnot(p \land q)$ i $\lnot p \lor \lnot q$ nisu jednake zato \v sto prva ima 6~simbola, a druga~5.
S druge strane, zbog tautologije
$$
  \vDash \lnot(p \land q) \Leftrightarrow (\lnot p \lor \lnot q)
$$
ove dve formule su \enquote{logi\v cki jednake} u smislu da je prva ta\v cna ako i samo ako je druga
ta\v cna. \v Cinjenicu da su dve formule \enquote{logi\v cki jednake} zato ozna\v cavamo novim simbolom:
$$
  \lnot(p \land q) \equiv (\lnot p \lor \lnot q)
$$
koga formalno defini\v semo ovako.

\begin{DEF}
  Iskazne formule $\phi$ i $\psi$ su \emph{ekvivalentne}, i pi\v semo $\phi \equiv \psi$, ako je
  $\vDash \phi \Leftrightarrow \psi$.
\end{DEF}

Simbol $\equiv$ ne treba me\v sati sa logi\v ckim veznikom $\Leftrightarrow$. Logi\v cki veznik
$\Leftrightarrow$ je deo jezika iskaznog ra\v cuna i koristi se kao simbol prilikom formiranja iskaznih formula.
Simbol $\equiv$ pripada \emph{metajeziku}, dakle, jeziku koga koristimo da pri\v camo o iskaznim formulama. Zapis
$\phi \equiv \psi$ \emph{nije logi\v cka formula}, ve\'c samo kra\'ci na\v cin da se ka\v ze
\enquote{formule $\phi$ i $\psi$ su ekvivalentne}. S druge strane zapis $\phi \Leftrightarrow \psi$
je iskazna formula.

% Slede\'ca teorema pokazuje da ekvivalentne formule imaju istu logi\v cku vrednost u svim valuacijama.
% 
% \begin{THM}\label{ft03.thm.taut1}
%   Neka su $\phi$ i $\psi$ iskazne formule. Tada $\phi \equiv \psi$ ako i samo ako za svaku valuaciju
%   $v$ iskaznih slova koja se javljaju u obe formule va\v zi $v(\phi) = v(\psi)$. 
% \end{THM}
% \begin{proof}[Dokaz]
%   $(\Rightarrow)$
%   Prema definiciji, $\phi \equiv \psi$ zna\v ci da je
%   $\phi \Leftrightarrow \psi$ tautologija, pa za svaku valuaciju $v$ va\v zi $v(\phi \Leftrightarrow \psi) = \top$.
%   Prema tabelici koja defini\v se $\Leftrightarrow$ kao logi\v cku operaciju vidimo da je
%   $v(\phi \Leftrightarrow \psi) = \top$ ta\v cno u onim situacijama u kojima je $v(\phi) = v(\psi)$.
%   Dakle, za svaku valuaciju $v$ je $v(\phi) = v(\psi)$.
% 
%   $(\Leftarrow)$
%   Pretpostavimo, sada, da za svaku valuaciju $v$ va\v zi $v(\phi) = v(\psi)$. Tada
%   prema tabelici koja defini\v se $\Leftrightarrow$ kao logi\v cku operaciju vidimo da
%   za svaku valuaciju $v$ va\v zi i da je $v(\phi \Leftrightarrow \psi) = \top$. To, po definiciji,
%   zna\v ci da je $\phi \Leftrightarrow \psi$ tautologija, odnosno, da je $\phi \equiv \psi$.
% \end{proof}
% 
% Sada \v zelimo da opi\v semo kako se formule transformi\v su jedna u drugu. Prvo \'cemo
% uvesti jednu pogodnu oznaku i onda \'cemo dokazati da se primenom \enquote{ekvivalentnih transformacija}
% od formule dobija njoj ekvivalentna formula.
% 
% Neka su $\alpha$ i $\phi$ iskazne formule i neka je $p$ neko iskazno slovo koje se javlja u iskaznoj formuli $\phi$.
% Ozna\v cimo sa
% $$
%     \phi[\alpha / p]
% $$
% formulu u kojoj je svako pojavljivanje slova $p$ zamenjenom formulom $\alpha$.
% (Oznaku $\phi[\alpha / p]$ \v citamo: \enquote{formula $\phi$ u kojoj $\alpha$ menja $p$}.)
% 
% \begin{EX}
%   Neka je $\phi = (p \Rightarrow (q \Rightarrow p))$ i $\alpha = (p \land \lnot q)$. Tada je
%   $$
%     \phi[\alpha / p] = ((p \land \lnot q) \Rightarrow (q \Rightarrow (p \land \lnot q))).
%   $$
%   Tako\dj e,
%   $$
%     \phi[\phi / q] = (p \Rightarrow ((p \Rightarrow (q \Rightarrow p)) \Rightarrow p)).
%   $$
% \end{EX}
% 
% \begin{THM}\label{ft03.thm.ekviv-smena}
%   $(a)$
%   Neka su $\alpha$ i $\beta$ iskazne formule takve da je $\vDash \alpha \Leftrightarrow \beta$, neka je $\phi$
%   iskazna formula i neka je $p$ neko iskazno slovo koje se javlja u iskaznoj formuli $\phi$.
%   Tada je $\vDash \phi[\alpha/p] \Leftrightarrow \phi[\beta/p]$.
% 
%   $(b)$
%   Neka su $\phi$ i $\psi$ iskazne formule takve da je $\phi \equiv \psi$,
%   neka je $p$ neko iskazno slovo koje se javlja i u formuli $\phi$ i u formuli $\psi$, i neka je
%   $\alpha$ proizvoljna iskazna formula.
%   Tada je $\phi[\alpha/p] \equiv \psi[\alpha/p]$.
% \end{THM}
% \begin{proof}[Dokaz]
%   $(a)$
%   Prema Teoremi~\ref{ft03.thm.taut1} dovoljno je pokazati da za svaku valuaciju $v$ va\v zi
%   $$
%     v(\phi[\alpha/p]) = v(\phi[\beta/p]).
%   $$
%   No, ovo je jasno: tamo gde je u formuli $\phi$ stajalo slovo $p$ sada stoji formula $\alpha$ na levoj
%   strani, odnosno, formula $\beta$ na desnoj strani. Po\v sto je $\alpha \Leftrightarrow \beta$ tautologija,
%   prema Teoremi~\ref{ft03.thm.taut1} znamo da je $v(\alpha) = v(\beta)$ za svaku valuaciju $v$.
%   Kako ostala slova u formuli $\phi$ nismo menjali, mora biti $v(\phi[\alpha/p]) = v(\phi[\beta/p])$.
% 
%   $(b)$ Sli\v cno kao $(a)$.
% \end{proof}

\begin{EX}
  Dokazati da je $p \land q \Rightarrow r \;\;\equiv\;\; p \Rightarrow (q \Rightarrow r)$.
\end{EX}
\begin{proof}[Re\v senje]
  Krenu\'cemo od formule sa desne strane i primeniti niz ekvivalentnih transformacija da bismo dobili formulu sa leve strane:
  \begin{align*}
    p \Rightarrow (q \Rightarrow r)
    &\;\;\equiv\;\; \lnot p \lor (q \Rightarrow r) && \text{Zakon materijalne implikacije}\\
    &\;\;\equiv\;\; \lnot p \lor (\lnot q \lor r)  && \text{Zakon materijalne implikacije}\\
    &\;\;\equiv\;\; (\lnot p \lor \lnot q) \lor r  && \text{Asocijativnost disjunkcije}\\
    &\;\;\equiv\;\; \lnot(p \land q) \lor r        && \text{De Morganov zakon}\\
    &\;\;\equiv\;\; (p \land q) \Rightarrow r      && \text{Zakon materijalne implikacije }\qedhere
  \end{align*}
\end{proof}

\medskip

\v Cesto je potrebno na\'ci negaciju neke formule (primere iz programerske prakse \'cemo videti kasnije).
\emph{Negirati formulu $\phi$} zna\v ci krenuti od formule $\lnot \phi$ i primeniti niz ekvivalentnih
transformacija tako da na kraju dobijemo formulu u kojoj negacija napada samo iskazna slova.

\begin{EX}\label{ft03.ex.neg-implic}
  Negirati formulu $p \Rightarrow q$.
\end{EX}
\begin{proof}[Re\v senje]
  \begin{align*}
    \lnot(p \Rightarrow q)
    &\;\;\equiv\;\; \lnot(\lnot p \lor q)       && \text{Zakon materijalne implikacije}\\
    &\;\;\equiv\;\; \lnot\lnot p \land \lnot q  && \text{De Morganov zakon}\\
    &\;\;\equiv\;\; p \land \lnot q             && \text{Involutivnost negacije }\qedhere
  \end{align*}
\end{proof}

\clearpage

\begin{EX}\label{ft03.ex.neg-ekviv}
  Negirati formulu $p \Leftrightarrow q$.
\end{EX}
\begin{proof}[Re\v senje]
  \begin{align*}
    \lnot(p \Leftrightarrow q)
    &\;\;\equiv\;\; \lnot((p \Rightarrow q) \land (q \Rightarrow p))    && \text{Zakon materijalne ekvivalencije}\\
    &\;\;\equiv\;\; \lnot(p \Rightarrow q) \lor \lnot(q \Rightarrow p)  && \text{De Morganov zakon}\\
    &\;\;\equiv\;\; (p \land \lnot q) \lor (q \land \lnot p)            && \text{Primer~\ref{ft03.ex.neg-implic}}\qedhere
  \end{align*}
\end{proof}

\begin{EX}
  Negirati formulu $p \lor q \Leftrightarrow p \lor r$.
\end{EX}
\begin{proof}[Re\v senje]
  \begin{align*}
    \lnot&(p \lor q \Leftrightarrow p \lor r) \;\;\equiv\;\;\\
    &\;\;\equiv\;\; ((p \lor q) \land \lnot (p \lor r)) \lor ((p \lor r) \land \lnot(p \lor q)) && \text{Primer~\ref{ft03.ex.neg-ekviv}}\\
    &\;\;\equiv\;\; ((p \lor q) \land \lnot p \land \lnot r)) \lor ((p \lor r) \land \lnot p \land \lnot q)) && \text{De Morganov zakon}\\
    &\;\;\equiv\;\; (p \land \lnot p \land \lnot r) \lor (q \land \lnot p \land \lnot r) \lor\\
    &\text{\hphantom{$\;\;\;\;\equiv\;\;$}}\lor (p \land \lnot p \land \lnot q) \lor (r \land \lnot p \land \lnot q) && \text{Distributivnost}\\
    &\;\;\equiv\;\; (\bot \land \lnot r) \lor (q \land \lnot p \land \lnot r) \lor\\
    &\text{\hphantom{$\;\;\;\;\equiv\;\;$}}\lor (\bot \land \lnot q) \lor (r \land \lnot p \land \lnot q) && \text{Princip kontradikcije}\\
    &\;\;\equiv\;\; \bot \lor (q \land \lnot p \land \lnot r) \lor \bot \lor (r \land \lnot p \land \lnot q) && \text{Apsorptivni element}\\
    &\;\;\equiv\;\; (q \land \lnot p \land \lnot r) \lor (r \land \lnot p \land \lnot q) && \text{Neutralni element}\qedhere
  \end{align*}
\end{proof}

\ifBIOINF
\else

    \bigskip
    
    Ekvivalentne transformacije se koriste jo\v s i da se formule dovedu do posebnih \enquote{pravilnih oblika}
    koje zovemo \emph{normalne forme}. Kada se formula dovede u normalnu formu neke njene odlike postaju
    jasno uo\v cljive, \v sto \'cemo pokazati na primerima.
    
    \begin{DEF}
      Iskazna formula $\phi$ je u \emph{konjunktivnoj normalnoj formi} ako je
      $$
        \phi = (\alpha_1 \lor \alpha_2 \lor \ldots \lor \alpha_l) \land
               (\beta_1  \lor \beta_2  \lor \ldots \lor \beta_m)  \land \ldots \land
               (\theta_1 \lor \theta_2 \lor \ldots \lor \theta_n)
      $$
      gde je svaka od formula $\alpha_i$, $\beta_j$, $\gamma_k$ iskazno slovo ili negirano iskazno slovo.
      Iskazna formula $\phi$ je u \emph{disjunktivnoj normalnoj formi} ako je
      $$
        \phi = (\alpha_1 \land \alpha_2 \land \ldots \land \alpha_l) \lor
               (\beta_1  \land \beta_2  \land \ldots \land \beta_m)  \lor \ldots \lor
               (\theta_1 \land \theta_2 \land \ldots \land \theta_n)
      $$
      gde je svaka od formula $\alpha_i$, $\beta_j$, $\gamma_k$ iskazno slovo ili negirano iskazno slovo.
    \end{DEF}
    
    Slede\'ca teorema (koju ne\'cemo dokazati) garantuje da se svaka iskazna formula mo\v ze ekvivalentnim
    transformacijama dovesti i u konjunktivnu normalnu formu i u disjunktivnu normalnu formu.
    
    \begin{THM}
      Za svaku formulu $\phi$ postoji formula $\psi_\land$ u konjunktivnoj normalnoj formi
      takva da je $\phi \equiv \psi_\land$, i postoji formula $\psi_\lor$
      u disjunktivnoj normalnoj formi takva da je $\phi \equiv \psi_\lor$.
    \end{THM}
    
    \begin{EX}
      Odrediti konjunktivnu i disjunktivnu normalnu formu formule $p \Leftrightarrow q$.
    \end{EX}
    \begin{proof}[Re\v senje]
      Konjunktivnu normalnu formu dobijamo slede\'cim nizom ekvivalentnih transformacija:
      \begin{align*}
        p \Leftrightarrow q
        &\;\;\equiv\;\; (p \Rightarrow q) \land (q \Rightarrow p)         && \text{Zakon materijalne ekvivalencije}\\
        &\;\;\equiv\;\; (\lnot p \lor q) \land (\lnot q \lor p)           && \text{Zakon materijalne implikacije}\\
      \end{align*}
      Disjunktivnu normalnu formu dobijamo ako prethodni niz transformacija nastavimo ovako:
      \begin{align*}
        &\;\;\equiv\;\; (\lnot p \land \lnot q) \lor (\lnot p \land p) \lor (q \land \lnot q) \lor (q \land p)   && \text{Distributivnost}\\
        &\;\;\equiv\;\; (\lnot p \land \lnot q) \lor \bot \lor \bot \lor (q \land p)   && \text{Zakon kontradikcije}\\
        &\;\;\equiv\;\; (\lnot p \land \lnot q) \lor (q \land p)   && \text{Neutralni element}\qedhere
      \end{align*}
    \end{proof}
    
    \begin{NAP}
      \textit{Primetimo da se iz konjunktivne normalne forme disjunktivna normalna forma dobija pravolinijski,
      primenom distributivnosti (i \enquote{\v ci\v s\'cenjem sme\'ca} koje
      pri tom nastaje). Obrnuto, iz disjunktivne normalne forme se lako dobija konjunktivna primenom distributivnosti
      (i \enquote{\v ci\v s\'cenjem sme\'ca}).}
    \end{NAP}
    
    \enlargethispage{\baselineskip}
    \begin{EX}
      Pokazati da je $\vDash \lnot(p \Rightarrow q) \lor \lnot p \lor q$ tautologija svo\dj enjem formule na konjunktivnu normalnu formu.
    \end{EX}
    \begin{proof}[Re\v senje]
      \begin{align*}
        \lnot(p \Rightarrow q) \lor \lnot p \lor q
        &\;\;\equiv\;\; \lnot(\lnot p \lor q) \lor \lnot p \lor q          && \text{Zakon mat.\ implikacije}\\
        &\;\;\equiv\;\; (\lnot\lnot p \land \lnot q) \lor \lnot p \lor q   && \text{De Morganov zakon}\\
        &\;\;\equiv\;\; (p \land \lnot q) \lor \lnot p \lor q              && \text{Involutivnost negacije}\\
        &\;\;\equiv\;\; (p \lor \lnot p \lor q) \land (\lnot q \lor \lnot p \lor q) && \text{Distributivnost}\\
        &\;\;\equiv\;\; (p \lor \lnot p \lor q) \land (\lnot q \lor q \lor \lnot p) && \text{Komutativnost}\\
        &\;\;\equiv\;\; (\top \lor q) \land (\top \lor \lnot p) && \text{Zakon isklju\v cenja tre\'ceg}\\
        &\;\;\equiv\;\; \top \land \top && \text{Neutralni element}\\
        &\;\;\equiv\;\; \top && \text{Neutralni element}\qedhere
      \end{align*}
    \end{proof}
    
    Prethodni primeri sugeri\v su proceduru za svo\dj enje bilo koje formule na konjunktivnu ili disjunktivnu normalnu formu:
    \begin{enumerate}\itemsep -1pt
    \item eliminisati simbole $\Leftrightarrow$ koriste\'ci Zakon materijalne ekvivalencije;
    \item eliminisati simbole $\Rightarrow$ koriste\'ci Zakon materijalne implikacije;
    \item \enquote{gurnuti} sve simbole $\lnot$ neposredno do iskaznih slova koriste\'ci De Morganove zakone;
    \item eliminisati $\lnot \lnot$ koriste\'ci Involutivnost negacije;
    \item svesti formulu na \v zeljenu normalnu formu koriste\'ci Distributivnost;
    \item (pri tome \enquote{\v cistimo sme\'ce} koriste\'ci Zakon isklju\v cenja tre\'ceg, Apsorptivni i Neutralni element).
    \end{enumerate}
\fi


\section{Tautologije kao \enquote{zakoni mi\v sljenja}?}

Formalno rezonovanje se ne mo\v ze primeniti na svakodnevne situacije pre svega zato \v sto
samo dve logi\v cke vrednosti, \emph{ta\v cno} i \emph{neta\v cno}, \v cesto nisu dovoljne da
obuhvate \v citav spektar zna\v cenja koja ljudi pripisuju situacijama u kojima se nalaze. Na primer,
ako poku\v samo da tautologiju
$$
  \vDash \lnot\lnot\phi \Leftrightarrow \phi
$$
primenimo na re\v cenicu
\begin{center}
  \enquote{Nije mi neprijatno}
\end{center}
sledilo bi da izjava \enquote{Nije mi neprijatno} ima isto zna\v cenje kao izjava \enquote{Prijatno mi je},
\v sto naj\v ce\v s\'ce nije tako. Ako \v zelim da ka\v zem da mi je prijatno, re\'ci \'cu direktno:
\enquote{Prijatno mi je}. Me\dj utim, ako ka\v zem: \enquote{Nije mi neprijatno}, to ne mora da zna\v ci da se ose\'cam prijatno.

Krajem XIX veka, Luis Kerol, autor \v cuvene knjige \textit{Alisa u zemlji \v cuda},
je bio jedan od najglasnijih pobornika ideje da se strategije formalnog rezonovanja
\emph{ne mogu primenjivati na svakodnevne situacije}. Takvo njegovo vi\dj enje formalne logike je bilo
u o\v stroj suprotnosti sa op\v ste prihva\'cenim stavom da je Bulova
\emph{Matemati\v cka analiza logike} siguran prvi korak ka ostvarenju Lajbnicovog sna o
\enquote{univerzalnom jeziku nauke}. Da bi ilustrovao svoj stav Luis Kerol
je pravio glavolomke kojima je dovodio \v citaoce u razne apsurdne situacije i zahtevao od njih da
primene pravila formalnog rezonovanja. Sada \'cemo navesti tri takva primera.

\paragraph{Primer: Problem sa terijerima i kometama.}
Iz slede\'ce tri premise izvesti odgovaraju\'ci zaklju\v cak:
\begin{enumerate}
\item
  Nijedan terijer ne luta me\dj u zvezdama.
\item
  Ni\v sta \v sto ne luta me\dj u zvezdama nije kometa.
\item
  Samo terijeri imaju kovrd\v zav rep.
\end{enumerate}

\noindent
\textit{Re\v senje.}
Uvedimo slede\'ce oznake:
$$
  \begin{array}{r@{\;=\;}l@{\qquad}r@{\;=\;}l}
    T & \text{terijer}, & L & \text{luta me\dj u zvezdama},\\
    K & \text{kometa},  & R & \text{ima kovrd\v zav rep}.
  \end{array}
$$
Prethodne tri izjave mo\v zemo formalno da zapi\v semo ovako:
\begin{enumerate}
\item
  $T \Rightarrow \lnot L$ (ako je terijer onda ne luta me\dj u zvezdama);
\item
  $\lnot L \Rightarrow \lnot K$ (ako ne luta me\dj u zvezdama onda nije kometa); i
\item
  $R \Rightarrow T$ (ako ima kovrd\v zav rep onda je to terijer).
\end{enumerate}
Sada se lako uo\v cava niz implikacija:
$$
  R \underset{(3)}\Rightarrow T \underset{(1)}\Rightarrow \lnot L \underset{(2)}\Rightarrow \lnot K.
$$
Dakle, $R \Rightarrow \lnot K$. Imaju\'ci u vidu da je
$
  R \Rightarrow \lnot K \equiv K \Rightarrow \lnot R
$
zaklju\v cujemo, tako\dj e, da je $K \Rightarrow \lnot R$, odnosno, da
\emph{komete nemaju kovrd\v zav rep.}

\paragraph{Primer: Problem sa magarcima i bivolima.}
Iz slede\'cih premisa izvesti odgovaraju\'ci zaklju\v cak:
\begin{enumerate}
\item%1
  \v Zivotinje koje se ne ritaju nije lako uznemiriti.
\item%2
  Magarci nemaju rogove.
\item%3
  Bivo uvek mo\v ze nekoga da baci preko ograde.
\item%4
  Nijednu \v zivotinju koja se rita nije lako progutati.
\item%5
  \v Zivotinje bez rogova ne mogu nikoga da bace preko ograde.
\item%6
  Sve \v zivotinje je lako uznemiriti, osim bivola.
\end{enumerate}

\noindent
\textit{Re\v senje.}
Uvedimo slede\'ce oznake:
$$
  \begin{array}{r@{\;=\;}l@{\qquad}r@{\;=\;}l}
  M & \text{magarac},    & U & \text{lako se uznemiri},\\
  B & \text{bivo},       & O & \text{mo\v ze nekoga da baci preko ograde},\\
  R & \text{rita se},    & P & \text{lako ga je progutati},\\
  G & \text{ima rogove}.
  \end{array}
$$
Tada prethodne izjave mo\v zemo formalno da zapi\v semo ovako:
\begin{enumerate}
\item%1
  $\lnot R \Rightarrow \lnot U$
  (ako se ne rita onda ga nije lako uznemiriti);
\item%2
  $M \Rightarrow \lnot G$
  (ako je magarac onda nema rogove);
\item%3
  $B \Rightarrow O$
  (ako je bivo onda mo\v ze nekoga da baci preko ograde);
\item%4
  $R \Rightarrow \lnot P$
  (ako se rita onda ga nije lako progutati);
\item%5
  $\lnot G \Rightarrow \lnot O$
  (ako nema rogove onda ne mo\v ze nikoga da baci preko ograde);
\item%6
  $\lnot B \Rightarrow U$
  (ako nije bivo onda se lako uznemiri).
\end{enumerate}
Iz (2), (5) i kontrapozicije stava (3) koja glasi $\lnot O \Rightarrow \lnot B$, tim redom, zaklju\v cujemo
$$
  M \Rightarrow \lnot B,
$$
odnosno, da \emph{magarac nije bivo}, a onda iz tog zaklju\v cka, (6), kontrapozicije stava (1) i stava (4), tim redom,
zaklju\v cujemo
$$
  M \Rightarrow \lnot P,
$$
odnosno, da \emph{magarca nije lako progutati}.

\paragraph{Primer: Problem sa svinjama i balonima.}
Iz slede\'cih premisa izvesti odgovaraju\'ci zaklju\v cak:
\begin{enumerate}
\item%1
  Svi koji niti hodaju po u\v zetu niti jedu jeftine kifle su stari.
\item%2
  Svinjama koje vole da se kin\dj ure se iskazuje po\v stovanje.
\item%3
  Mudar pilot balona\footnote{ovde se misli na balone sa toplim vazduhom kojima mo\v ze da se putuje kao u knjizi \v Zila Verna
  \enquote{Put oko sveta za osamdeset dana}} uvek nosi ki\v sobran.
\item%4
  Niko ko je luckast i jede jeftine kifle ne treba da jede na javnom mestu.
\item%5
  Mlada stvorenja koja vole da pilotiraju balonom vole da se kin\dj ure.
\item%6
  Debela stvorenja koja su luckasta mogu da jedu na javnom mestu, ali samo ako ne hodaju po u\v zetu.
\item%7
  Nijedno mudro stvorenje ne hoda po u\v zetu, ako voli da se kin\dj uri.
\item%8
  Svinja koja nosi ki\v sobran je luckasta.
\item%9
  Svi koji ne hodaju po u\v zetu i kojima se iskazuje po\v stovanje su debeli.
\end{enumerate}

\noindent
\textit{Re\v senje.}
Uvedimo slede\'ce oznake:
\newcommand\NJ{\mathit{Nj}}%
$$
  \begin{array}{r@{\;=\;}l@{\qquad}r@{\;=\;}l}
  D   & \text{debeo},              & V    & \text{voli da se kin\dj uri},\\
  \NJ & \text{svinja},             & T    & \text{mo\v ze da jede na javnom mestu},\\
  S   & \text{star},               & P    & \text{iskazuje mu se po\v stovanje},\\
  M   & \text{mudar},              & F    & \text{jede jeftine kifle},\\
  L   & \text{luckast},            & B    & \text{pilotira balonom},\\
  K   & \text{nosi ki\v sobran},   & U    & \text{hoda po u\v zetu}.\\
  \end{array}
$$
Tada prethodne izjave mo\v zemo formalno da zapi\v semo ovako:
\begin{enumerate}
\item%1
  $\lnot U \land \lnot F \Rightarrow S$
  (ako ne hoda po u\v zetu i ne jede jeftine kifle onda je star);
\item%2
  $\NJ \land V \Rightarrow P$
  (ako je svinja i voli da se kin\dj uri onda mu se iskazuje po\v stovanje);
\item%3
  $M \land B \Rightarrow K$
  (ako je mudar i pilotira balonom onda nosi ki\v sobran);
\item%4
  $L \land F \Rightarrow \lnot T$
  (ako je luckast i jede jeftine kifle onda ne treba da jede na javnom mestu);
\item%5
  $\lnot S \land B \Rightarrow V$
  (ako je mlad ($=$ nije star) i pilotira balonom onda voli da se kin\dj uri);
\item%6
  $D \land L \land \lnot U \Rightarrow T$
  (ako je debeo i luckast, i jo\v s ako ne hoda po u\v zetu, onda mo\v ze da jede na javnom mestu);
\item%7
  $M \land V \Rightarrow \lnot U$
  (ako je mudar i voli da se kin\dj uri onda ne hoda po u\v zetu);
\item%8
  $\NJ \land K \Rightarrow L$
  (ako je svinja i nosi ki\v sobran onda je luckast);
\item%9
  $\lnot U \land P \Rightarrow D$
  (ako ne hoda po u\v zetu i iskazuje mu se po\v stovanje onda je debeo).
\end{enumerate}
Za razliku od prethodna dva primera gde smo imali jednostavne nizove implikacija, ovde je proces zaklju\v civanja slo\v zeniji
zato \v sto sa leve strane svake implikacije imamo konjunkciju uslova.

Interesantno je da je za ovu vrstu iskaza Luis Kerol razvio posebnu deduktivnu proceduru koju je detaljno opisao u
svojoj knjizi \enquote{Simboli\v cka logika}.\footnote{Lewis Carroll, \textit{Symbolic Logic}, MACMILLAN AND CO., Limited,
New York: The MacMillan Company, 1897}
Ova procedura je ponovo otkrivena 1960-ih godina i danas predstavlja osnovu sistema za automatsko rezonovanje
koji se zove \emph{rezolucija}.

Metod koga Luis Kerol predla\v ze veoma li\v ci na re\v savanje sistema linearnih jedna\v cina Gausovim metodom eliminacije.
Uo\v cimo prvo da je
$$
  \phi \Rightarrow \psi \quad\equiv\quad \phi \land \lnot \psi \Rightarrow \bot,
$$
\v sto se lako pokazuje. Ovo veoma li\v ci na slede\'cu algebarsku transformaciju:
$$
  p = q \quad\text{ako i samo ako}\quad p + (-q) = 0.
$$
Prvi korak u re\v savanju problema se sastoji u tome da, koriste\'ci navedeno pravilo,
implikacije (1)--(9) zapi\v semo u obliku \enquote{homogenog sistema jedna\v cina}:
\begin{enumerate}
\item%1
  $\lnot U \land \lnot F \land \lnot S \Rightarrow \bot$
\item%2
  $\NJ \land V \land \lnot P \Rightarrow \bot$
\item%3
  $M \land B \land \lnot K \Rightarrow \bot$
\item%4
  $L \land F \land T \Rightarrow \bot$
\item%5
  $\lnot S \land B \land \lnot V \Rightarrow \bot$
\item%6
  $D \land L \land \lnot U \land \lnot T \Rightarrow \bot$
\item%7
  $M \land V \land U \Rightarrow \bot$
\item%8
  $\NJ \land K \land \lnot L \Rightarrow \bot$
\item%9
  $\lnot U \land P \land \lnot D \Rightarrow \bot$.
\end{enumerate}
Ovaj \enquote{homogeni sistem jedna\v cina} se mo\v ze \enquote{re\v siti} metodom \enquote{eliminacije slova}. Naime, lako se vidi da:
$$
  \text{iz\quad} a \land c \Rightarrow \bot \text{\quad i\quad} b \land \lnot c \Rightarrow \bot \text{\quad sledi\quad}
  a \land b \Rightarrow \bot,
$$
(slovo $c$ je eliminisano), \v sto direktno odgovara ovom algebarskom tvr\dj enju:
$$
  \text{iz\quad} a + c = 0 \text{\quad i\quad} b + (- c) = 0 \text{\quad sledi\quad}
  a + b = 0.
$$
Metodom eliminacije slova dobijamo slede\'ci niz zaklju\v caka, koriste\'ci vi\v se puta da je $\phi \land \phi \equiv \phi$:
\begin{enumerate}\setcounter{enumi}{9}
\item%10
  $\lnot U \land \lnot S \land L \land T \Rightarrow \bot$
  (eliminacijom slova $F$ iz (1) i (4))
\item%11
  $\lnot U \land \lnot S \land L \land D \Rightarrow \bot$
  (eliminacijom slova $T$ iz (10) i (6))
\item%12
  $\lnot U \land \lnot S \land D \land \NJ \land K \Rightarrow \bot$
  (eliminacijom slova $L$ iz (11) i (8))
\item%13
  $\lnot U \land \lnot S \land \NJ \land K \land P \Rightarrow \bot$
  (eliminacijom slova $D$ iz (12) i (9))
\item%14
  $\lnot S \land \NJ \land K \land P \land M \land V \Rightarrow \bot$
  (eliminacijom slova $U$ iz (13) i (7))
\item%15
  $\lnot S \land \NJ \land P \land M \land V \land B \Rightarrow \bot$
  (eliminacijom slova $K$ iz (14) i (3))
\item%16
  $\lnot S \land \NJ \land M \land V \land B \Rightarrow \bot$
  (eliminacijom slova $P$ iz (15) i (2))
\item%17
  $\lnot S \land \NJ \land M \land B \Rightarrow \bot$
  (eliminacijom slova $V$ iz (16) i (5))
\end{enumerate}
Ukoliko zaklju\v cak (17) zapi\v semo u ekvivalentnom obliku
$$
  \lnot S \land \NJ \land B \Rightarrow \lnot M
$$
i ukoliko, kao i ranije, $\lnot S$ zamenimo sa \enquote{mlad},
kao zaklju\v cak dobijamo da \emph{mlade svinje koje pilotiraju balonom nisu mudre}.



%%%%%%%%%%%%%%%%%%%%%%%%%%%%%%%%%%%%%%%%%%%%%%%%%%%%%%%%%%%%%%%%%%%%%%%%%%%%%%
% \section{Iskazni ra\v cun kao formalna teorija}
% 
% Poka\v zimo sada da formalne teorije mogu da opi\v su i fenomene koji odgovaraju
% (formalnom) logi\v ckom rezonovanju. Kao primer \'cemo pokazati formalnu teoriju
% koja opisuje iskazni ra\v cun u slede\'cem smislu: niz simbola je posledica ove
% formalne teorije ako i samo ako je odgovaraju\'ca formula tautologija.
% 
% Jezik ove formalne teorije je jezik iskaznog ra\v cuna: \v cine ga iskazna slova
% $p_1$, $p_2$,~\ldots, logi\v cki veznici $\lnot$ i $\Rightarrow$, i zagrade ( i ),
% a formule su sve formule iskaznog ra\v cuna u kojima se javljaju samo ova dva logi\v cka veznika.
% Va\v zno je re\'ci da odluka da se fokusiramo samo na logi\v cke veznike $\lnot$ i $\Rightarrow$ ne predstavlja
% nikakvo ograni\v cenje u izra\v zajnoj mo\'ci teorije, zato \v sto se ostali veznici mogu u ovako
% osiroma\v senu teoriju uvesti pomo\'cu slede\'cih tautologija:
% 
% \medskip
% 
% \begin{tabular}{l}
%   $\vDash p \lor q \Leftrightarrow (\lnot p \Rightarrow q)$ \\
%   $\vDash p \land q \Leftrightarrow \lnot(p \Rightarrow \lnot q)$ \\
%   $\vDash (p \Leftrightarrow q) \Leftrightarrow (p \Rightarrow q) \land (q \Rightarrow p).$
% \end{tabular}
% 
% \medskip
% 
% \noindent
% Prema tome, u formalnoj teoriji iskaznog ra\v cuna preostale logi\v cke veznike shvatamo
% kao skra\'ceni zapis nekih formula:
% \begin{itemize}\itemsep -2pt
% \item $p \lor q$ je skra\'ceni zapis za $\lnot p \Rightarrow q$,
% \item $p \land q$ je skra\'ceni zapis za $\lnot(p \Rightarrow \lnot q)$, i
% \item $p \Leftrightarrow q$ je skra\'ceni zapis za $(p \Rightarrow q) \land (q \Rightarrow p)$.
% \end{itemize}
% 
% \noindent
% U verziji koju \'cemo ovde pokazati formalna teorija iskaznog ra\v cuna ima \v cetiri aksiome:
% 
% \begin{description}
% \item[A1:] $p_1 \Rightarrow (p_2 \Rightarrow p_1)$;
% \item[A2:] $(p_1 \Rightarrow (p_2 \Rightarrow p_3)) \Rightarrow ((p_1 \Rightarrow p_2) \Rightarrow (p_1 \Rightarrow p_3))$;
% \item[A3:] $(\lnot p_2 \Rightarrow \lnot p_1) \Rightarrow (p_1 \Rightarrow p_2)$;
% \item[A4:] $\lnot p_1 \Rightarrow (p_1 \Rightarrow p_2)$.
% \end{description}
% i dva pravila izvo\dj enja --- \emph{Modus ponens} i \emph{Pravilo supstitucije}:
% $$
%   \mathit{MP}: \begin{array}{c}
%     \phi \qquad \phi \Rightarrow \psi \\ \hline
%     \psi
%   \end{array},\qquad
%   \mathit{SUP}: \begin{array}{cl}
%     \phi & \text{(koja sadr\v zi slovo $p$)} \\ \cline{1-1}
%     \phi[\alpha / p] & \text{(gde je $\alpha$ \emph{proizvoljna} formula).}
%   \end{array}
% $$
% 
% \medskip
% 
% Jasno je da nije svaka formula iskaznog ra\v cuna posledica ove formalne teorije.
% Ako je $\phi$ posledica ove formalne teorije pi\v semo $\vdash \phi$.
% 
% \begin{EX}\label{ft03.ex.alpha=>alpha}
%   Dokazati da je $\vdash \alpha \Rightarrow \alpha$ za svaku iskaznu formulu $\alpha$.
% \end{EX}
% \begin{proof}[Re\v senje]
%   Neka je $\alpha$ proizvolja formula i neka je $x$ iskazno slovo koje se ne
%   javlja u formuli~$\alpha$ i koje je razli\v cito od~$p_1$, $p_2$ i $p_3$. Tada slede\'ci
%   niz predstavlja dokaz formule $\alpha \Rightarrow \alpha$ u formalnoj teoriji iskaznog ra\v cuna:
% \begin{enumerate}
% \item
%   $p_1 \Rightarrow (p_2 \Rightarrow p_1)$\newline
%   [A1]
% \item
%   $p_1 \Rightarrow (x \Rightarrow p_1)$\newline
%   [SUP na (1) gde $x$ menja $p_2$]
% \item
%   $\alpha \Rightarrow (x \Rightarrow \alpha)$\newline
%   [SUP na (2) gde $\alpha$ menja $p_1$]
% \item
%   $\alpha \Rightarrow ((\alpha \Rightarrow \alpha) \Rightarrow \alpha)$\newline
%   [SUP na (3) gde $(\alpha \Rightarrow \alpha)$ menja $x$]
% \item
%   $(p_1 \Rightarrow (p_2 \Rightarrow p_3)) \Rightarrow ((p_1 \Rightarrow p_2) \Rightarrow (p_1 \Rightarrow p_3))$\newline
%   [A2]
% \item
%   $(p_1 \Rightarrow (x \Rightarrow p_3)) \Rightarrow ((p_1 \Rightarrow x) \Rightarrow (p_1 \Rightarrow p_3))$\newline
%   [SUP na (5) gde $x$ menja $p_2$]
% \item
%   $(p_1 \Rightarrow (x \Rightarrow p_1)) \Rightarrow ((p_1 \Rightarrow x) \Rightarrow (p_1 \Rightarrow p_1))$\newline
%   [SUP na (6) gde $p_1$ menja $p_3$]
% \item
%   $(\alpha \Rightarrow (x \Rightarrow \alpha)) \Rightarrow ((\alpha \Rightarrow x) \Rightarrow (\alpha \Rightarrow \alpha))$\newline
%   [SUP na (7) gde $\alpha$ menja $p_1$]
% \item
%   $(\alpha \Rightarrow ((\alpha \Rightarrow \alpha) \Rightarrow \alpha)) \Rightarrow ((\alpha \Rightarrow (\alpha \Rightarrow \alpha)) \Rightarrow (\alpha \Rightarrow \alpha))$\newline
%   [SUP na (8) gde $(\alpha \Rightarrow \alpha)$ menja $x$]
% \item
%   $(\alpha \Rightarrow (\alpha \Rightarrow \alpha)) \Rightarrow (\alpha \Rightarrow \alpha)$\newline
%   [MP na (4) i (9)]
% \item
%   $\alpha \Rightarrow (\alpha \Rightarrow \alpha)$\newline
%   [SUP na (3) gde $\alpha$ menja $x$]
% \item
%   $\alpha \Rightarrow \alpha$\newline
%   [MP na (10) i (11)]\qedhere
% \end{enumerate}
% \end{proof}
% 
% \medskip
% 
% Ve\'c ovaj jednostavan primer pokazuje da \'ce rad u formalnoj teoriji iskaznog ra\v cuna biti jedno mukotrpno
% iskustvo. Zato \'cemo sada (bez dokaza) dati dve korisne teoreme koje u mnogome olak\v savaju
% na\v se (formalne) \v zivote.
% 
% Neka su $\alpha_1$, \ldots, $\alpha_n$ i $\phi$ iskazne formule i neka se u formuli $\phi$ javljaju
% slova $p_1$, \ldots, $p_n$. Ozna\v cimo sa
% $$
%     \phi[\alpha_1 / p_1, \ldots, \alpha_n / p_n]
% $$
% formulu u kojoj se istovremeno svako pojavljivanje slova $p_i$ zameni formulom $\alpha_i$, $1 \le i \le n$.
% 
% \begin{EX}
%   Neka je $\phi = p_1 \Rightarrow (p_2 \Rightarrow (p_1 \Rightarrow p_2))$ i neka je
%   $\alpha_1 = p_1 \Rightarrow p_2$, a $\alpha_2 = p_2 \Rightarrow p_1$. Tada je
%   $$
%     \phi[\alpha_1 / p_1, \alpha_2 / p_2] =
%     (p_1 \Rightarrow p_2) \Rightarrow ((p_2 \Rightarrow p_1) \Rightarrow ((p_1 \Rightarrow p_2) \Rightarrow (p_2 \Rightarrow p_1))).
%   $$
% \end{EX}
% 
% \begin{THM}[Teorema o istovremenoj supstituciji]
%   Neka su $\alpha_1$, \ldots, $\alpha_n$ i $\phi$ iskazne formule i neka se u formuli $\phi$ javljaju
%   slova $p_1$, \ldots, $p_n$. Ako je $\vdash \phi$ onda je
%   $\vdash \phi[\alpha_1 / p_1, \ldots, \alpha_n / p_n]$.
% \end{THM}
% 
% Druga va\v zna teorema koju \'cemo predstaviti je Teorema dedukcije za iskazni ra\v cun, za \v ciju
% formulaciju nam treba jedan va\v zan pojam.
% 
% \begin{DEF}
%   Neka je $\Sigma$ skup formula iskaznog ra\v cuna i neka je $\phi$ neka formula iskaznog ra\v cuna. Tada oznaka
%   $\Sigma \vdash \phi$ zna\v ci da postoji kona\v can niz formula
%   $$
%     \alpha_1, \; \alpha_2, \; \ldots, \; \alpha_n
%   $$
%   takav da je $\alpha_n = \phi$, i za svako $i \in \{1, \ldots, n\}$ formula $\alpha_i$ je ili aksioma, ili
%   neka od formula iz $\Sigma$, ili je dobijena primenom nekog pravila izvo\dj enja na neke od formula koje joj
%   prethode u nizu. Navedeni niz formula tada nazivamo \emph{dokaz formule $\phi$ iz pretpostavki $\Sigma$}.
% \end{DEF}
% 
% \begin{THM}[Teorema dedukcije za iskazni ra\v cun]
%   Neka je $\Sigma$ skup formula iskaznog ra\v cuna i neka su $\phi$ i $\psi$ formule iskaznog ra\v cuna. Tada
%   $\Sigma \vdash \phi \Rightarrow \psi$ ako i samo ako $\Sigma \union \{\phi\} \vdash \psi$.
% \end{THM}
% 
% \begin{EX}
%   Dokazati da je $\vdash \alpha \Rightarrow \alpha$ za svaku iskaznu formulu $\alpha$.
% \end{EX}
% \begin{proof}[Re\v senje]
%   Prema Teoremi dedukcije dovoljno je dokazati da je
%   $$
%     \alpha \vdash \alpha
%   $$
%   \v sto sada mo\v zemo da uradimo u jednom koraku:
%   \begin{enumerate}
%   \item
%     $\alpha$\newline
%     [hipoteza]
%   \end{enumerate}
%   Uporediti ovaj dokaz sa dokazom koji je dat u Primeru~\ref{ft03.ex.alpha=>alpha}.
% \end{proof}
% 
% \begin{EX}\label{ft03.ex.tranz}
%   Dokazati da je $\vdash (p \Rightarrow q) \Rightarrow ((q \Rightarrow r) \Rightarrow (p \Rightarrow r))$.
% \end{EX}
% \begin{proof}[Re\v senje]
%   Prema Teoremi dedukcije dovoljno je dokazati da je
%   $$
%     p \Rightarrow q \vdash (q \Rightarrow r) \Rightarrow (p \Rightarrow r).
%   $$
%   Ako jo\v s jednom primenimo Teoremu dedukcije dobijamo da treba pokazati da je
%   $$
%     p \Rightarrow q, q \Rightarrow r \vdash p \Rightarrow r.
%   $$
%   Nakon tre\'ce primene Teoreme dedukcije dobijamo da je, zapravo, dovoljno pokazati
%   $$
%     p \Rightarrow q, q \Rightarrow r, p \vdash r
%   $$
%   \v sto je veoma lako:
% \begin{enumerate}
% \item
%   $p$\newline
%   [hipoteza]
% \item
%   $p \Rightarrow q$\newline
%   [hipoteza]
% \item
%   $q$\newline
%   [MP na (1) i (2)]
% \item
%   $q \Rightarrow r$\newline
%   [hipoteza]
% \item
%   $r$\newline
%   [MP na (3) i (4)]\qedhere
% \end{enumerate}
% \end{proof}
% 
% \begin{EX}\label{ft03.ex.lnot-lnot-p--p}
%   Dokazati da je $\vdash \lnot\lnot p \Rightarrow p$.
% \end{EX}
% \begin{proof}[Re\v senje]
%   Prema Teoremi dedukcije dovoljno je dokazati da je $\lnot\lnot p \vdash p$.
% \begin{enumerate}
% \item
%   $\lnot\lnot p$\newline
%   [hipoteza]
% \item
%   $\lnot\lnot p \Rightarrow (\lnot p \Rightarrow \lnot\lnot\lnot p)$\newline
%   [Teorema istovr.\ supstitucije na A4 gde $\lnot p / p_1$ i $\lnot\lnot\lnot p / p_2$]
% \item
%   $\lnot p \Rightarrow \lnot\lnot\lnot p$\newline
%   [MP na (1) i (2)]
% \item
%   $(\lnot p \Rightarrow \lnot\lnot\lnot p) \Rightarrow (\lnot\lnot p \Rightarrow p)$\newline
%   [Teorema istovr.\  supstitucije na A3 gde $\lnot\lnot p / p_1$ i $p / p_2$]
% \item
%   $\lnot\lnot p \Rightarrow p$\newline
%   [MP na (3) i (4)]
% \item
%   $p$\newline
%   [MP na (1) i (5)]\qedhere
% \end{enumerate}
% \end{proof}
% 
% \begin{EX}\label{ft03.ex.p--lnot-lnot-p}
%   Dokazati da je $\vdash p \Rightarrow \lnot\lnot p$.
% \end{EX}
% \begin{proof}[Re\v senje]
%   U Primeru~\ref{ft03.ex.lnot-lnot-p--p} smo dokazali da je $\vdash \lnot\lnot p \Rightarrow p$.
%   To zna\v ci da postoji kona\v can niz formula koji predstavlja dokaz ove formule.
%   Taj kona\v can niz formula mo\v zemo umetnuti gdegod nam zatreba i tako unutar nekog drugog dokaza
%   dobiti, kao me\dj ukorak, dokaz ove formule. Odatle sledi da sve \v sto smo ranije dokazali
%   mo\v zemo da koristimo u dokazima novih formula. To \'cemo sada iskoristiti.
% \begin{enumerate}
% \item
%   $(\lnot\lnot\lnot p \Rightarrow \lnot p) \Rightarrow (p \Rightarrow \lnot\lnot p)$\newline
%   [Teorema istovr.\ supstitucije na A3 gde $\lnot\lnot p / p_2$ i $p / p_1$]
% \item
%   $\lnot\lnot p \Rightarrow p$\newline
%   [Primer~\ref{ft03.ex.lnot-lnot-p--p}]
% \item
%   $\lnot\lnot\lnot p \Rightarrow \lnot p$\newline
%   [SUP na (2) gde $\lnot p$ menja $p$]
% \item
%   $p \Rightarrow \lnot\lnot p$\newline
%   [MP na (3) i (1)]\qedhere
% \end{enumerate}
% \end{proof}
% 
% \begin{EX}
%   Dokazati da je $\vdash (p \Rightarrow q) \Rightarrow (\lnot q \Rightarrow \lnot p)$.
% \end{EX}
% \begin{proof}[Re\v senje]
%   Prema Teoremi dedukcije dovoljno je pokazati da $p \Rightarrow q \vdash \lnot q \Rightarrow \lnot p$.
% \begin{enumerate}
% \item
%   $p \Rightarrow q$\newline
%   [hipoteza]
% \item
%   $q \Rightarrow \lnot\lnot q$\newline
%   [SUP na Primer~\ref{ft03.ex.p--lnot-lnot-p} gde $q$ menja $p$]
% \item
%   $(p \Rightarrow q) \Rightarrow ((q \Rightarrow \lnot\lnot q) \Rightarrow (p \Rightarrow \lnot\lnot q))$\newline
%   [SUP na Primer~\ref{ft03.ex.tranz} gde $\lnot\lnot q$ menja $r$]
% \item
%   $(q \Rightarrow \lnot\lnot q) \Rightarrow (p \Rightarrow \lnot\lnot q)$\newline
%   [MP na (1) i (3)]
% \item
%   $p \Rightarrow \lnot\lnot q$\newline
%   [MP na (2) i (4)]
% \item
%   $\lnot\lnot p \Rightarrow p$\newline
%   [Primer~\ref{ft03.ex.lnot-lnot-p--p}]
% \item
%   $(\lnot\lnot p \Rightarrow p) \Rightarrow ((p \Rightarrow \lnot\lnot q) \Rightarrow (\lnot\lnot p \Rightarrow \lnot\lnot q))$\newline
%   [Teorema istovr.\ supstitucije na Primer~\ref{ft03.ex.tranz} gde $\lnot\lnot p / p$, $p/q$, $\lnot\lnot q / r$]
% \item
%   $(p \Rightarrow \lnot\lnot q) \Rightarrow (\lnot\lnot p \Rightarrow \lnot\lnot q)$\newline
%   [MP na (6) i (7)]
% \item
%   $\lnot\lnot p \Rightarrow \lnot\lnot q$\newline
%   [MP na (5) i (8)]
% \item
%   $(\lnot\lnot p \Rightarrow \lnot\lnot q) \Rightarrow (\lnot q \Rightarrow \lnot p)$\newline
%   [Teorema istovr.\ supstitucije na A3 gde $\lnot q / p_1$, $\lnot p / p_2$]
% \item
%   $\lnot q \Rightarrow \lnot p$\newline
%   [MP na (9) i (10)]\qedhere
% \end{enumerate}
% \end{proof}
% 
% \medskip
% 
% Ovo je samo nekoliko primera koji pokazuju kako se za neke tautologije pokazuje da su posledice
% formalne teorije iskaznog ra\v cuna. Uz mnogo truda mo\v ze se pokazati da su sve tautologije
% posledice formalne teorije iskaznog ra\v cuna, i nijedna druga formula to nije.
% Dakle, tautologije mo\v zemo precizno opisati formalnim teorijama.
% 
% \begin{THM}
%   Neka je $\phi$ iskazna formula. Tada $\vdash \phi$ ako i samo ako je $\phi$ tautologija.
% \end{THM}
% 
% Ako izjavu $\vDash \phi$ shvatimo kao \enquote{$\phi$ je ta\v cno}, a izjavu $\vdash \phi$ kao
% \enquote{$\phi$ je dokazivo} onda prethodna veoma va\v zna teorema mo\v ze da se shvati i ovako:
% formula iskaznog ra\v cuna je ta\v cna ako i samo ako je dokaziva.
% 
% Ovo je zna\v cajna osobina jer pokazuje da se u iskaznom ra\v cunu \emph{semantika} (pojam ta\v cnosti)
% i \emph{sintaksa} (pojam formalne dokazivosti) poklapaju, \v sto nije uvek slu\v caj!
% 
% %%%%%%%%%%%%%%%%%%%%%%%%%%%%%%%%%%%%%%%%%%%%%%%%%%%%%%%%%%%%%%%%%%%%%%%%%%%%%%%%%%%%%



\section{Iskazne formule u analizi programskog koda}

Iskazne formule u programiranju odgovaraju uslovima koji se javljaju u \enquote{if}, \enquote{while} i sli\v cnim
kontrolnom strukturama. Sada \'cemo na nekoliko jednostavnih primera
pokazati kako se ve\v stine koje smo stekli koriste za analizu toka programskog koda,
detekciju mrtvog koda i analizu petlji. Kako ra\v cunarski programi postaju sve slo\v zeniji
i slo\v zeniji, formalna analiza programskog koda postaje va\v zna ve\v stina za jednog programera
jer, kao \v sto je 1972.\ rekao Edgar Dajkstra,
\begin{quote}
  \enquote{Program testing can be used to show the presence of bugs, but never their absence.}
  
  Edsgar Dijkstra
\end{quote}

\clearpage

\begin{EX}
  Posmatrajmo slede\'ci metod:
{\small
\begin{verbatim}
    static int NZD(int a, int b) {
      if(a <= 0 || b <= 0) { return -1; }
      // (1)
      ...
    }
\end{verbatim}
}
\noindent
\v Sta va\v zi u ta\v cki (1) koda?
\end{EX}
\begin{proof}[Re\v senje]
  Ako metod stigne u ta\v cku (1) to zna\v ci da uslov iza \enquote{if} nije ispunjen. Dakle, u ta\v cki (1)
  va\v zi
  $$
    \lnot(a \le 0 \lor b \le 0)
  $$
  \v sto je, prema De Morganovom zakonu, ekvivalentno sa
  $$
    a > 0 \land b > 0.
  $$
  Dakle, ako metod pro\dj e \enquote{if} kontrolnu strukturu i nastavi sa radom znamo da je $a > 0$ i $b > 0$.
\end{proof}

\begin{EX}
  U slede\'cem programskom fragmentu
{\small
\begin{verbatim}
    if(x >= 6 && y < 7) {
      // (1)
      ...
    }
    else {
      // (2)
      ...
    }
\end{verbatim}
}
\end{EX}
\noindent
\v sta va\v zi u ta\v cki (1) koda, a \v sta u ta\v cki (2)?
\begin{proof}[Re\v senje]
  Ako kod stigne u ta\v cku (1) tada je uslov u \enquote{if} ispunjen, pa u ta\v cki (1) va\v zi
  $x \ge 6 \land y < 7$.
  S druge strane, ako je kod stigao do ta\v cke (2) uslov u \enquote{if} nije ispunjen, tako da u ta\v cki (2) va\v zi
  $\lnot(x \ge 6 \land y < 7)$,
  \v sto je ekvivalentno sa $x < 6 \lor y \ge 7$.
\end{proof}

\begin{EX}
  Koji uslovi va\v ze u ta\v ckama (1), (2), (3) i (4) slede\'ceg programskog fragmenta:
{\small
\begin{verbatim}
  if(x < 1) {
    // (1)
    ...
  }
  else if(x < 5) {
    // (2)
    ...
  }
  else if(x > 10) {
    // (3)
    ...
  }
  else {
    // (4)
    ...
  }
\end{verbatim}
}
\end{EX}
\begin{proof}[Re\v senje]
  U ta\v cki (1) va\v zi uslov naveden u prvom \enquote{if}, dakle $x < 1$.
  
  Ako je kod stigao u ta\v cku (2), zna\v ci da uslov
  iza prvog \enquote{if} nije ispunjen, ali uslov iza drugog jeste. Zato u ta\v cki (2) va\v zi
  $\lnot(x < 1) \land x < 5$ \v sto je ekvivalentno sa $x \ge 1 \land x < 5$,
  odnosno, $1 \le x < 5$.

  Ako je kod stigao u ta\v cku (3), zna\v ci da uslovi
  iza prva dva \enquote{if}-a nisu ispunjeni, ali uslov iza tre\'ceg jeste. Zato u ta\v cki (3) va\v zi
  $\lnot(x < 1) \land \lnot(x < 5) \land x > 10$, \v sto je ekvivalentno sa
  $x \ge 1 \land x \ge 5 \land x > 10$, odnosno, $x > 10$.
  
  Kona\v cno, ako je kod stigao u ta\v cku (4) nijedan od tri uslova nije ispunjen. Zato u ta\v cki (4)
  va\v zi $\lnot(x < 1) \land \lnot(x < 5) \land \lnot(x > 10)$, \v sto je ekvivalentno sa
  $x \ge 1 \land x \ge 5 \land x \le 10$, odnosno, $5 \le x \le 10$.
\end{proof}

\begin{EX}
  Koji uslovi va\v ze u ta\v ckama (1), (2), (3) i (4) slede\'ceg programskog fragmenta:
{\small
\begin{verbatim}
  if(x < 1) {
    // (1)
    ...
  }
  else if(y > 5) {
    // (2)
    ...
  }
  else if(x > 3) {
    // (3)
    ...
  }
  else {
    // (4)
    ...
  }
\end{verbatim}
}
\end{EX}

\begin{proof}[Re\v senje]
  U ta\v cki (1) va\v zi $x < 1$.
  
  U ta\v cki (2) va\v zi $\lnot(x < 1) \land y > 5$ \v sto je ekvivalentno sa $x \ge 1 \land y > 5$.

  U ta\v cki (3) va\v zi $\lnot(x < 1) \land \lnot(y > 5) \land x > 3$, \v sto je ekvivalentno sa
  $x \ge 1 \land y \le 5 \land x > 3$, odnosno, $x > 3 \land y \le 5$.
  
  Kona\v cno, u ta\v cki (4)
  va\v zi $\lnot(x < 1) \land \lnot(y > 5) \land \lnot(x > 3)$, \v sto je ekvivalentno sa
  $x \ge 1 \land y \le 5 \land x \le 3$, odnosno, $1 \le x \le 3 \land y \le 5$.
\end{proof}


\begin{EX}
  Analiza uslova se mo\v ze upotrebiti da se detektuje mrtav kod (\v sto je kod koji se nikada
  ne\'ce izvr\v siti). U slede\'cem programskom fragmentu detektovati mrtav kod:
{\small
\begin{verbatim}
  if(x < 1) {
    // (1)
    ...
  }
  else if(x < 5) {
    // (2)
    ...
  }
  else if(x < -3) {
    // (3)
    ...
  }
  else {
    // (4)
    ...
  }
\end{verbatim}
}
\end{EX}
\begin{proof}[Re\v senje]
  U ta\v cki (1) va\v zi uslov naveden u prvom \enquote{if}, dakle $x < 1$.
  
  U ta\v cki (2) va\v zi $\lnot(x < 1) \land x < 5$ \v sto je ekvivalentno sa $x \ge 1 \land x < 5$,
  odnosno, $1 \le x < 5$.
  
  U ta\v cki (3) va\v zi
  $\lnot(x < 1) \land \lnot(x < 5) \land x < -3$, \v sto je ekvivalentno sa
  $x \ge 1 \land x \ge 5 \land x < -3$, \v sto nije mogu\'ce! Dakle, kod naveden u bloku koji po\v cinje u
  ta\v cki (3) je mrtav kod.
  
  Kona\v cno, u ta\v cki (4)
  va\v zi $\lnot(x < 1) \land \lnot(x < 5) \land \lnot(x < -3)$, \v sto je ekvivalentno sa
  $x \ge 1 \land x \ge 5 \land x \ge -3$, odnosno, $x \ge 5$.
\end{proof}

\begin{EX}\label{ft03.ex.swap}
  \v Sta va\v zi u ta\v ckama (1), (2), (3) i (4) slede\'ceg programskog fragmenta?
  \v Sta radi ovaj fragment?
\end{EX}
{\small
\begin{verbatim}
  // ucitamo celobrojne promenljive a i b
  int x = a; int y = b;
  // (1)
  int pom = x;
  // (2)
\end{verbatim}
\clearpage
\begin{verbatim}
  x = y;
  // (3)
  y = pom;
  // (4)
\end{verbatim}
}
\begin{proof}[Re\v senje]
  U ta\v cki (1) va\v zi
  $$
    x = a \land y = b.
  $$
  U ta\v cki (2) va\v zi
  $$
    x = a \land y = b \land \mathit{pom} = a.
  $$
  U ta\v cki (3) va\v zi
  $$
    x = b \land y = b \land \mathit{pom} = a,
  $$
  a u ta\v cki (4) va\v zi
  $$
    x = b \land y = a \land \mathit{pom} = a.
  $$
  Dakle, ovaj fragment razmenjuje sadr\v zaj promenljivih \verb"x" i \verb"y".
\end{proof}

\begin{EX}
  \v Sta va\v zi u ta\v ckama (1), (2), (3) i (4) slede\'ceg programskog fragmenta:

\begin{tabular}{l}
  \verb"// ucitamo celobrojne promenljive a i b"\\
  \verb"int x = a; int y = b;"\\
  \verb"// (1)"\\
  \verb"x = x + y;"\\
  \verb"// (2)"\\
  \verb"y = x - y;"\\
  \verb"// (3)"\\
  \verb"x = x - y;"\\
  \verb"// (4)"
\end{tabular}

\noindent
\v Sta radi ovaj fragment?
\end{EX}
\begin{proof}[Re\v senje]
  U ta\v cki (1) va\v zi
  $$
    x = a \land y = b.
  $$
  Odatle je
  $$
    x + y = a + b \land y = b,
  $$
  pa u ta\v cki (2) va\v zi
  $$
    x = a + b \land y = b.
  $$
  Odavde sledi da je
  $$
    x = a + b \land x - y = a,
  $$
  pa u ta\v cki (3) va\v zi
  $$
    x = a + b \land y = a.
  $$
  Odavde sledi da je
  $$
    x - y = b \land y = a,
  $$
  pa u ta\v cki (4) va\v zi
  $$
    x = b \land y = a.
  $$
  Dakle, ovaj fragment razmenjuje sadr\v zaj promenljivih \verb"x" i \verb"y". Ovaj kod je samo akademska glavolomka!
  \emph{Nikada ne treba na ovaj na\v cin razmenjivati vrednosti promenljivih!}
  Smislen na\v cin je onaj pokazan u Primeru~\ref{ft03.ex.swap}.
\end{proof}

\begin{EX}
  Neka je $\phi$ nekakav logi\v cki izraz.
  Koji uslovi va\v ze u ta\v ckama (1) i (2) slede\'ceg programskog fragmenta:

\medskip

\begin{tabular}{l}
  \verb"while("$\phi$\verb") {"\\
  \verb"  // (1)"\\
  \verb"  ..."\\
  \verb"}"\\
  \verb"// (2)"
\end{tabular}
\end{EX}
\begin{proof}[Re\v senje]
  Ta\v cka (1) se nalazi odmah nakon ulaska u petlju, tako da tu sigurno va\v zi $\phi$.
  Ta\v cka (2) se nalazi neposredno nakon izlaska iz petlje, pa tu sigurno va\v zi $\lnot\phi$ jer se telo
  petlje izvr\v sava sve dok je $\phi$ ta\v cno.
\end{proof}

\begin{EX}
  Neka je $\phi$ nekakav logi\v cki izraz.
  Koji uslovi va\v ze u ta\v ckama (1) i (2) slede\'ceg programskog fragmenta:

\medskip

\begin{tabular}{l}
  \verb"do {"\\
  \verb"  // (1)"\\
  \verb"  ..."\\
  \verb"} while("$\phi$\verb");"\\
  \verb"// (2)"
\end{tabular}
\end{EX}
\begin{proof}[Re\v senje]
  Prilikom prvog ulaska u petlju nemamo dovoljno informacija da bismo znali \v sta va\v zi u ta\v cki (1).
  No, prilikom svakog narednog ciklusa petlje u ta\v cki (1) sigurno va\v zi $\phi$.
  Ta\v cka (2) se nalazi neposredno nakon izlaska iz petlje, pa tu sigurno va\v zi $\lnot\phi$ jer se telo
  petlje izvr\v sava sve dok je $\phi$ ta\v cno.
\end{proof}


\begin{EX}
  Posmatrajmo slede\'ci metod:
{\small
\begin{verbatim}
    static int MN(int x, int y) {
      if(x <= 0 || y <= 0) { return -1; }
      // (1)
      int p = 0;
      int u = x;
      // (2)
      do {
        // (3)
        p += y;
        u--;
        // (4)
      } while(u > 0);
      // (5)
      return p;
    }
\end{verbatim}
}

\medskip

\noindent
$(a)$ \v Sta va\v zi u ta\v ckama (1) i (2) koda?

\medskip

\noindent
$(b)$ Dokazati da u ta\v cki (3) koda va\v zi $p + u\cdot y = x \cdot y \land u > 0$, a da u ta\v cki (4) va\v zi $p + u\cdot y = x \cdot y \land u \ge 0$.

\medskip

\noindent
$(c)$ \v Sta va\v zi u ta\v cki (5) koda? \v Sta radi metod \verb"MN"?
\end{EX}
\begin{proof}[Re\v senje]
  $(a)$
  Ako metod stigne u ta\v cku (1) to zna\v ci da uslov iza \enquote{if} nije ispunjen. Dakle, u ta\v cki (1)
  va\v zi
  $$
    x > 0 \land y > 0.
  $$
  Onda u ta\v cki (2) va\v zi
  $$
    x > 0 \land y > 0 \land p = 0 \land u = x.
  $$

  $(b)$
  Kada program prvi put u\dj e u petlju u ta\v cki (3) va\v zi
  $$
    x > 0 \land y > 0 \land p = 0 \land u = x
  $$
  odakle lako sledi da je
  $$
    p + u\cdot y = x \cdot y \land u > 0.
  $$
  Ovo je ekvivalentno sa
  $$
    (p + y) + (u - 1)\cdot y = x \cdot y \land u - 1 \ge 0
  $$
  zato \v sto je $u$ celobrojna promenljiva. Kada nakon naredbi
  {\small
  \begin{verbatim}
    p += y;
    u--;
  \end{verbatim}
  }
  \noindent
  stignemo u ta\v cku (4) promenljiva $p$ se uve\'ca za $y$, promenljiva $u$ se smanji za 1 i zato u ta\v cki (4) va\v zi:
  $$
    p + u\cdot y = x \cdot y \land u \ge 0
  $$
  (jer $p + y$ postaje novo $p$, a $u - 1$ postaje novo $u$). Ako ciklus nastavi sa radom, ponovo sti\v zemo u ta\v cku (3)
  u kojoj sada va\v zi
  $$
    (p + u\cdot y = x \cdot y \land u \ge 0) \land (u > 0)
  $$
  \v sto je ekvivalentno sa
  $$
    p + u\cdot y = x \cdot y \land u > 0.
  $$
  Nakon toga na isti na\v cin kao ranije zaklju\v cujemo da u ta\v cki (4) va\v zi
  $$
    p + u\cdot y = x \cdot y \land u \ge 0.
  $$
  Dakle, tokom rada petlje sve vreme va\v zi
  $$
    p + u\cdot y = x \cdot y \land u \ge 0.
  $$
  Ova formula se zato zove \emph{invarijanta petlje} jer je invarijantna (ne menja oblik) tokom izvr\v savanja petlje.
  
  $(c)$
  U ta\v cki (5) koda va\v zi
  $$
    (p + u\cdot y = x \cdot y \land u \ge 0) \land \lnot(u > 0)
  $$
  odnosno,
  $$
    p + u\cdot y = x \cdot y \land (u \ge 0 \land \lnot(u > 0)).
  $$
  Kako je $u \ge 0 \land \lnot (u > 0)$ ekvivalentno sa $u = 0$ dobijamo da u ta\v cki (5) zapravo va\v zi:
  $$
    p + 0\cdot y = x \cdot y \text{\quad tj.\quad} p = x\cdot y.
  $$
  Prema tome, metod \verb"MN" radi slede\'ce: ako je $x \le 0$ ili $y \le 0$ metod vra\'ca $-1$, dok u ostalim slu\v cajevima
  ra\v cuna proizvod $x \cdot y$.
\end{proof}

U prethodnom primeru smo videli kako se invarijante petlje koriste da analiziramo rad petlje. \emph{Invarijanta petlje} je formula
koja va\v zi pre ulaska u telo petlje i nakon svakog izvr\v savanja tela petlje. Onda ne menja svoj oblik tokom rada petlje, odatle ime.
Drugi va\v zan aspekt analize rada petlji je dokaz da \'ce se petlja zavr\v siti nakon kona\v cnog broja koraka. To nije uvek
o\v cigledno. Zato \'cemo u slede\'cem primeru pokazati jednu standardnu tehniku za ovakve dokaze. Ideja je da se uo\v ci
algebarski izraz koji se zove \emph{varijanta petlje} koji se u svakom prolasku kroz telo petlje smanjuje i petlja se zavr\v sava
kada vrednost ovog izraza stigne do neke unapred zadate vrednosti.

\begin{EX}
  Dokazati da se metod \verb"MN" uvek zavr\v sava nakon kona\v cnog broja koraka:
{\small
\begin{verbatim}
    static int MN(int x, int y) {
      if(x <= 0 || y <= 0) { return -1; }
      int p = 0;
      int u = x;
\end{verbatim}
\clearpage
\begin{verbatim}
      do {
        p += y;
        u--;
      } while(u > 0);
      return p;
    }
\end{verbatim}
}
\end{EX}
\begin{proof}[Re\v senje]
  Lako se vidi da promenljiva $u$ slu\v zi kao varijanta petlje: ona na po\v cetku dobija neku pozitivnu vrednost,
  u svakom prolazu kroz telo petlje se njena vrednost smanjuje za 1, i petlja se zavr\v sava kada $u$ postane~0.
\end{proof}

\begin{EX}
  Posmatrajmo slede\'ci metod:
{\small
\begin{verbatim}
    static int CK(int x, int y) {
      if(x <= 0 || y <= 0) { return -1; }
      int q = 0;
      int r = x;
      // (1)
      while(r >= y) {
        // (2)
        r -= y;
        q++;
        // (3)
      }
      // (4)
      return q;
    }
\end{verbatim}
}

\medskip

\noindent
$(a)$ \v Sta va\v zi u ta\v cki (1)?

\medskip

\noindent
$(b)$ Dokazati da u ta\v cki (2) koda va\v zi $q \cdot y + r = x \land r \ge y$,
a da u ta\v cki (3) va\v zi $q \cdot y + r = x \land r \ge 0$ (ovo je invarijanta petlje).

\medskip

\noindent
$(c)$ \v Sta va\v zi u ta\v cki (4) koda? \v Sta radi metod \verb"CK"?

\medskip

\noindent
$(d)$ Dokazati da se ovaj metod uvek zavr\v sava.
\end{EX}
\begin{proof}[Re\v senje]
  $(a)$
  Ako metod stigne u ta\v cku (1) to zna\v ci da uslov iza \enquote{if} nije ispunjen. Dakle, u ta\v cki (1)
  va\v zi
  $$
    q = 0 \land r = x \land x > 0 \land y > 0.
  $$

  $(b)$
  Kada program prvi put u\dj e u petlju u ta\v cki (2) va\v zi
  $$
    q = 0 \land r = x \land x > 0 \land y > 0 \land r \ge y
  $$
  odakle lako sledi da je
  $$
    q \cdot y + r = x \land r \ge y.
  $$
  Ovo je ekvivalentno sa
  $$
    (q + 1) \cdot y + (r - y) = x \land r - y \ge 0.
  $$
  Kada nakon naredbi
  {\small
  \begin{verbatim}
      r -= y;
      q++;
  \end{verbatim}
  }
  \noindent
  stignemo u ta\v cku (3) promenljiva $q$ se uve\'ca $1$, promenljiva $r$ se smanji za $y$ i zato u ta\v cki (3) va\v zi:
  $$
    q \cdot y + r = x \land r \ge 0
  $$
  (jer $q + 1$ postaje novo $q$, a $r - y$ postaje novo $r$). Ciklus nastavlja sa radom, pa ponovo sti\v zemo u ta\v cku (2)
  u kojoj sada va\v zi
  $$
    q \cdot y + r = x \land r \ge y.
  $$
  Nakon toga na isti na\v cin kao ranije zaklju\v cujemo da u ta\v cki (3) va\v zi
  $$
    p + u\cdot y = x \cdot y \land u \ge 0.
  $$
  Dakle, tokom rada petlje sve vreme va\v zi
  $$
    q \cdot y + r = x \land r \ge 0.
  $$
  Dalje ne moramo da se vrtimo u krug, sve je jasno.
  
  $(c)$
  U ta\v cki (4) koda va\v zi
  $$
    (q \cdot y + r = x \land r \ge 0) \land \lnot(r \ge y)
  $$
  odnosno,
  $$
    q \cdot y + r = x \land 0 \le r < y.
  $$
  Prema tome, metod \verb"CK" radi slede\'ce: ako je $x \le 0$ ili $y \le 0$ metod vra\'ca $-1$, dok u ostalim slu\v cajevima
  ra\v cuna koli\v cnik i ostatak pri celobrojnom deljenju $x$ sa $y$, i kao rezultat vra\'ca celibrojni koli\v cnik.

  $(d)$
  Lako se vidi da promenljiva $r$ slu\v zi kao varijanta petlje: ona na po\v cetku dobija neku pozitivnu vrednost,
  u svakom prolazu kroz telo petlje se njena vrednost smanjuje za $y$, i petlja se zavr\v sava kada $r$ postane manje od~$y$.
\end{proof}

\clearpage

\begin{EX}
  Posmatrajmo slede\'ci metod:
{\small
\begin{verbatim}
    static int NZD(int x, int y) {
      if(x <= 0 || y <= 0) { return -1; }
      int a = x;
      int b = y;
      int r;
      // (1)
      do {
        // (2)
        r = a % b;
        a = b;
        b = r;
        // (3)
      } while(b > 0);
      // (4)
      return a;
    }
\end{verbatim}
}

\medskip

\noindent
$(a)$ \v Sta va\v zi u ta\v cki (1)?

\medskip

\noindent
$(b)$ Dokazati da u ta\v cki (2) koda va\v zi $\NZD(a, b) = \NZD(x, y) \land a > 0 \land b > 0$
(gde smo sa $\NZD(x, y)$ ozna\v cili najve\'ci zajedni\v cki delilac brojeva $x$ i $y$),
a da u ta\v cki (3) va\v zi $\NZD(a, b) = \NZD(x, y) \land a > 0 \land b \ge 0$ (ovo je invarijanta petlje).

\medskip

\noindent
$(c)$ Dokazati da u ta\v cki (4) koda va\v zi $a = \NZD(x, y)$.

\medskip

\noindent
$(d)$ Dokazati da se ovaj metod uvek zavr\v sava.
\end{EX}
\begin{proof}[Re\v senje]
  $(a)$
  Ako metod stigne u ta\v cku (1) to zna\v ci da uslov iza \enquote{if} nije ispunjen. Dakle, u ta\v cki (1)
  va\v zi
  $$
    a = x > 0 \land b = y > 0.
  $$

  $(b)$
  Kada program prvi put u\dj e u petlju u ta\v cki (2) va\v zi
  $$
    a = x > 0 \land b = y > 0
  $$
  odakle je jasno da je
  $$
    \NZD(a, b) = \NZD(x, y) \land a > 0 \land b > 0.
  $$
  Obzirom da je
  $$
    \NZD(a, b) = \NZD(b, a \% b)
  $$
  (elementarna teorija brojeva), u ta\v cki (2), prema tome, va\v zi:
  $$
    \NZD(b, a \% b) = \NZD(x, y) \land a > 0 \land b > 0.
  $$
  Evo kako se ova formula transformi\v se dok se izvr\v savaju tri naredbe dodele u telu ciklusa:
  \begin{center}
    \begin{tabular}{ll}
      Formula: & $\NZD(b, a \% b) = \NZD(x, y) \land a > 0 \land b > 0$\\
      Naredba: & \verb"r = a % b;"\\
      Formula: & $\NZD(b, r) = \NZD(x, y) \land a > 0 \land b > 0 \land r \ge 0$\\
      Naredba: & \verb"a = b;"\\
      Formula: & $\NZD(a, r) = \NZD(x, y) \land a > 0 \land r \ge 0$\\
      Naredba: & \verb"b = r;"\\
      Formula: & $\NZD(a, b) = \NZD(x, y) \land a > 0 \land b \ge 0$.
    \end{tabular}
  \end{center}
  Dakle, u ta\v cki (3) va\v zi
  $$
    \NZD(a, b) = \NZD(x, y) \land a > 0 \land b \ge 0.
  $$
  Ako ciklus nastavi sa radom, ponovo sti\v zemo u ta\v cku (2)
  u kojoj sada va\v zi
  $$
    \NZD(a, b) = \NZD(x, y) \land a > 0 \land b > 0.
  $$

  $(c)$
  U ta\v cki (4) koda va\v zi invarijanta petlje, i uslov za izlazak iz petlje:
  $$
    (\NZD(a, b) = \NZD(x, y) \land a > 0 \land b \ge 0) \land (b \le 0),
  $$
  odnosno,
  $$
    \NZD(a, b) = \NZD(x, y) \land a > 0 \land b = 0.
  $$
  Kako je $\NZD(x, y) = \NZD(a, b) = \NZD(a, 0) = a$ zaklju\v cujemo da
  metod \verb"NZD" radi slede\'ce: ako je $x \le 0$ ili $y \le 0$ metod vra\'ca $-1$, dok u ostalim slu\v cajevima
  odre\dj uje $\NZD(x, y)$.

  $(d)$
  Lako se vidi da promenljiva $b$ slu\v zi kao varijanta petlje: ona na po\v cetku dobija neku pozitivnu vrednost,
  u svakom prolazu kroz telo petlje se njena vrednost smanjuje (jer nova vrednost promenljive $b$ postaje ostatak pri celobrojnom deljenju
  $a$ sa $b$, \v sto je uvek manje od stare vrednost promenljive $b$), i petlja se zavr\v sava kada $b$ padne na~$0$.
\end{proof}


\section{Zadaci}

\begin{NAP}\it
Za prirodni broj $k$ i cele brojeve $m$ i $n$ pi\v semo $n\equiv_k m$ ako je $n-m$ deljivo brojem $k$. Tada ka\v zemo da je su brojevi $n$ i $m$ kongruentni po modulu $k$.
U matemati\v ckoj literaturi se umesto $n\equiv_k m$ \v ce\v s\'ce koristi oznaka $n \equiv m \pmod k$, ali \'ce
za na\v se potrebe skra\'ceni oblik $\equiv_k$ biti korisniji.
\end{NAP}

\clearpage

\ZAD{01.zad.1}{Proveriti da li su slede\' ci iskazi ta\v cni ili neta\v cni:
\begin{multicols}{2}
\begin{enumerate}[(a)]
\item $1\equiv_3 4$;
\item $4\equiv_7 22$;
\item $13\equiv_6 25$;
\item $2\equiv_5 -2$;
\item $3\equiv_7 -4$;
\item $5\equiv_{10} -5$;
\item $1701\equiv_9 1404$;
\item $2070\equiv_{11} 2205$.
\end{enumerate}
\end{multicols}
}

\ZAD{01.zad.2}{
Dati su slede\' ci iskazi:
\begin{align*}
p:&\text{ Mi\' ca se duri.}\\
q:&\text{ Joci je danas ro\dj endan.} \\
r:&\text{ Peca kupuje poklon.}
\end{align*}
Napisati slede\' ce iskaze re\v cima:
\begin{multicols}{3}
\begin{enumerate}[(a)]
\item $\neg p\wedge q$;
\item $p\vee q$;
\item $\neg p\Rightarrow q$;
\item $q\Leftrightarrow p$;
\item $q\Rightarrow r\wedge\neg p$;
\item $(q\Rightarrow r)\wedge p$.
\end{enumerate}
\end{multicols}
}

\ZAD{01.zad.3}{
Neka su $p$, $q$ i $r$ slede\' ci iskazi:
\begin{align*}
p:&\text{ \v Si\v smi\v si su slepi.}\\
q:&\text{ Bubamare jedu va\v si.}\\
r:&\text{ Mravi imaju duga\v cke zube.} 
\end{align*}
Zapisati slede\' ce re\v cenice pomo\' cu iskaznih formula:
\begin{enumerate}[(a)]
\item Ako su \v si\v smi\v si slepi, onda bubamare ne jedu va\v si.
\item \v Si\v smi\v si su slepi ako i samo ako mravi nemaju duge zube.
\item Mravi nemaju duge zube, i ako su \v si\v simi\v si slepi, onda bubamare ne jedu va\v si.
\item \v Si\v smi\v si su slepi ili bubamare jedu va\v si, i mravi nemaju duge zube ako bubamare ne jedu va\v si.
\end{enumerate}
}

\ZAD{01.zad.4}{
Negirati iskaze:
\begin{enumerate}[(a)]
\item \enquote{Sini\v sa je iz Novog Sada.}
\item \enquote{Kamion je bele boje.}
\item \enquote{Kombi je skrenuo desno.}
\item \enquote{Pera i Vlada su plavi.}
\item \enquote{Svi studenti su znali odgovor na pitanje.}
\item \enquote{Ako je Nikola visok, on igra ko\v sarku.}
\end{enumerate}
}

\begin{NAP}\it
Iskazna formula je \emph{kontradikcija} ako je neta\v cna za svaku valuaciju slova koja se u njoj javljaju.
Drugim re\v cima, kontradikcija je negacija tautologije.
\end{NAP}

\ZAD{01.zad.5}{Proveriti za svaku od  slede\' cih iskaznih formula da li je tautologija, kontradikcija, ili ni\v sta od ta dva:
\begin{enumerate}[(a)]
\item $p\wedge (p\vee q)\Leftrightarrow p$;
\item $\neg(p\wedge q)\Leftrightarrow(\neg p\vee\neg q)$;
\item $(\neg q\Rightarrow\neg p)\Rightarrow(p\Rightarrow q)$;
\item $p\wedge(p\Rightarrow q)\Rightarrow q$;
\item $(p\wedge q)\wedge(p\vee\neg q)$;
\item $p\wedge(q\vee r)\Leftrightarrow (p\wedge q)\vee(p\wedge r)$;
\item $p\Rightarrow p\vee q$;
\item $(p\Rightarrow q)\Leftrightarrow (\neg p\vee q)$;
\item $\big(p\Rightarrow(q\wedge r)\big)\Rightarrow\big((p\Rightarrow q)\wedge(p\Rightarrow r)\big)$;
\item $(p\Rightarrow(q\Rightarrow r))\Leftrightarrow (p\wedge \neg q)\vee r$;
\item $\big((p\Rightarrow\neg q)\vee r\big)\Leftrightarrow\big((q\Rightarrow r)\vee\neg p\big)$.
\end{enumerate}
}

\begin{NAP}\it
U kontekstu ovog ud\v zbenika, \emph{negirati iskaznu formulu $\phi$} zna\v ci zapisati formulu $\lnot \phi$
i onda je, koriste\'ci pogodne tautologije, svesti na ekvivalentan oblik kod koga logi\v cki veznik~$\lnot$ napada samo iskazna slova.
\end{NAP}

\ZAD{02.zad.1}{
Sastaviti tablicu istinitosnih vrednosti za slede\' ce iskazne formule, proveriti da li je ona tautologija, i negirati one iskazne formule koje nisu tautologije:
\begin{enumerate}[(a)]
\item $(p\Rightarrow(\neg q\Rightarrow r))\Leftrightarrow((p\wedge\neg q)\Rightarrow r)$;
\item $(p\Rightarrow(q\Rightarrow \neg r))\Leftrightarrow(r\Rightarrow (\neg p\vee \neg q))$;
\item $(\neg p\Rightarrow(q\Rightarrow r))\Leftrightarrow(p\vee\neg q\vee r)$;
\item $(\neg p\vee(q\Rightarrow  r))\Leftrightarrow(\neg r\Rightarrow \neg(p\wedge  q))$;
\item $(\neg p\vee  q)\Rightarrow(\neg p\wedge \neg q\wedge r)$;
\item $(\neg p\vee  \neg r)\Rightarrow(\neg p\wedge q\wedge r)$;
\item $(\neg p\vee(q\Rightarrow  r))\Leftrightarrow(\neg r\Rightarrow (p\Rightarrow\neg  q))$;
\item $(p\vee  \neg q\vee \neg r)\Rightarrow(\neg p\wedge \neg q\wedge r)$;
\item $(\neg p\wedge  \neg q)\Rightarrow(p\vee\neg q\vee \neg r)$.
\end{enumerate}
}

\ZAD{02.zad.2}{
Bez sastavljanja tablice istinitosnih vrednosti dokazati da su slede\' ce iskazne formule tautologije:
\begin{enumerate}[(a)]
\item $(p\Leftrightarrow \neg q)\wedge(\neg r\vee \neg s\Rightarrow q)\Rightarrow(p\Rightarrow s)$;
\item $s\wedge(\neg s\vee p)\Rightarrow(q\vee r\Rightarrow p)\vee t$;
\item $s\wedge p\wedge (s\vee \neg r\Rightarrow (\neg q\vee t))\Rightarrow (q\Rightarrow t)$;
\item $p\wedge (q\Rightarrow \neg p)\Rightarrow \neg q\vee s\vee \neg t$;
\item $s\wedge t\wedge(t\vee r\Rightarrow \neg p)\Rightarrow (q\Rightarrow \neg p)$;
\item $(t\Rightarrow s\vee \neg p)\wedge \neg s\wedge (p\wedge r)\Rightarrow (q\Rightarrow \neg t)$;
\item$(p\vee\neg r\Rightarrow \neg(s\vee t))\wedge(q\Rightarrow s)\Rightarrow r\vee \neg q$.
\end{enumerate}
}

